% Bab ini telah diubah pada 1/12/05.
%\setcounter{chapter}{2}
\chapter{Kombinatorika}\label{chp 3} 

\section{Permutasi}\label{sec 3.1}

Banyak masalah dalam teori peluang mengharuskan kita menghitung banyaknya cara suatu
kejadian tertentu dapat terjadi.  Untuk itu, kita mempelajari topik {\em permutasi} dan
{\em kombinasi.}  Pada bagian ini kita membahas permutasi, sedangkan kombinasi dibahas
pada bagian berikutnya.

Sebelum membahas permutasi, ada baiknya kita memperkenalkan suatu teknik pencacahan umum
yang akan memungkinkan kita menyelesaikan berbagai masalah pencacahan, termasuk masalah
menghitung banyaknya permutasi yang mungkin dari $n$ objek.
 
\subsection*{Masalah Pencacahan}

Perhatikan suatu percobaan yang berlangsung dalam beberapa tahap sedemikian rupa sehingga
banyaknya hasil $m$ pada tahap ke-$n$ tidak bergantung pada hasil tahap-tahap sebelumnya.
Bilangan $m$ dapat berbeda untuk tahap yang berbeda.  Kita ingin menghitung banyaknya cara
seluruh percobaan tersebut dapat dilakukan.

\begin{example}\label{exam 3.1} Anda sedang makan di restoran \'Emile's\index{Emile's
restaurant}, dan pramusaji memberi tahu bahwa Anda memiliki (a)~dua pilihan makanan
pembuka: sup atau jus;  (b)~tiga pilihan hidangan utama: hidangan daging, ikan, atau
sayuran; dan (c)~dua pilihan hidangan penutup: es krim atau kue.  Berapa banyak pilihan
yang mungkin untuk hidangan lengkap Anda?  Kita menggambarkan hidangan-hidangan yang
mungkin dengan diagram pohon pada Gambar~\ref{fig 3.1}.  Menu Anda ditentukan dalam tiga
tahap---pada setiap tahap, banyaknya pilihan yang mungkin tidak bergantung pada pilihan
di tahap-tahap sebelumnya: dua pilihan pada tahap pertama, tiga pada tahap kedua, dan
dua pada tahap ketiga.  Dari diagram pohon terlihat bahwa jumlah seluruh pilihan adalah
hasil kali banyaknya pilihan pada setiap tahap.  Dalam contoh ini terdapat
$2 \cdot 3 \cdot 2 = 12$ menu yang mungkin.  Contoh menu kita merupakan contoh teknik
pencacahan umum berikut.
\end{example}

\putfig{4truein}{PSfig3-1}{Pohon untuk menu Anda.}{fig 3.1} 

\subsection*{Teknik Pencacahan}

Suatu tugas akan dilakukan dalam urutan $r$ tahap.  Terdapat $n_1$ cara untuk
melakukan tahap pertama; untuk setiap dari $n_1$ cara tersebut, terdapat $n_2$ cara
untuk melakukan tahap kedua; untuk setiap dari $n_2$ cara tersebut, terdapat
$n_3$ cara untuk melakukan tahap ketiga, dan seterusnya.  Dengan demikian, jumlah
seluruh cara untuk menyelesaikan tugas itu diberikan oleh hasil kali
$N = n_1 \cdot n_2 \cdot \dots \cdot n_r$.

\putfig{4truein}{PSfig3-2}{Penetapan peluang dua tahap.}{fig 3.2} 
 
\subsection*{Diagram Pohon}\index{diagram pohon}

Diagram pohon sering berguna ketika kita mempelajari peluang kejadian yang berkaitan
dengan percobaan yang berlangsung dalam beberapa tahap dan peluang hasil pada setiap
tahap telah diberikan.  Sebagai contoh, anggaplah pemilik restoran \'Emile's mengamati
bahwa 80 persen pelanggannya memilih sup sebagai makanan pembuka dan 20 persen memilih
jus.  Di antara mereka yang memilih sup, 50 persen memilih daging, 30 persen memilih
ikan, dan 20 persen memilih hidangan sayuran.  Di antara mereka yang memilih jus sebagai
makanan pembuka, 30 persen memilih daging, 40 persen memilih ikan, dan 30 persen memilih
hidangan sayuran.  Kita dapat menggunakan informasi ini untuk memperkirakan peluang pada
dua tahap pertama sebagaimana ditunjukkan pada diagram pohon di Gambar~\ref{fig 3.2}.

Sebagai ruang sampel, kita memilih himpunan $\Omega$ yang berisi semua lintasan yang
mungkin, $\omega =
\omega_1$,~$\omega_2$, \dots,~$\omega_6$, melalui pohon tersebut.
Bagaimanakah kita harus menetapkan distribusi peluangnya?  Sebagai contoh, peluang apa
yang harus kita tetapkan bagi pelanggan yang memilih sup lalu daging?  Jika 8/10
pelanggan memilih sup dan kemudian 1/2 di antaranya memilih daging, proporsi pelanggan
yang memilih sup lalu daging adalah $8/10 \cdot 1/2 = 4/10$.  Hal ini menyarankan agar
distribusi peluang untuk setiap lintasan melalui pohon dipilih sebagai {\em hasil kali}
peluang pada setiap tahap di sepanjang lintasan tersebut.  Hasilnya adalah distribusi
peluang untuk titik-titik sampel $\omega$ yang ditunjukkan pada Gambar~\ref{fig 3.2}.
(Perhatikan bahwa $m(\omega_1) + \cdots +
m(\omega_6) = 1$.)  Dari sini terlihat,
misalnya, bahwa peluang seorang pelanggan memilih daging adalah
$m(\omega_1) + m(\omega_4) = .46$.
\par
Kita akan membahas ukuran pada pohon ini lebih lanjut ketika membicarakan konsep
peluang bersyarat dalam Bab~\ref{chp 4}.    Sekarang kita kembali ke masalah pencacahan
lainnya.

\begin{example}\label{exam 3.2} Kita dapat menunjukkan bahwa setidaknya ada dua orang di
Columbus, Ohio, yang memiliki tiga inisial yang sama.  Dengan menganggap setiap orang
memiliki tiga inisial, terdapat 26 kemungkinan untuk inisial pertama, 26 untuk inisial
kedua, dan 26 untuk inisial ketiga.  Jadi, terdapat $26^3 = 17{,}576$ susunan inisial
yang mungkin.  Angka ini lebih kecil daripada jumlah penduduk Columbus, Ohio; karena
itu, pasti ada setidaknya dua orang dengan tiga inisial yang sama.
\end{example}

Selanjutnya kita membahas masalah ulang tahun yang terkenal---masalah yang sering
digunakan untuk menunjukkan bahwa intuisi naif tidak selalu dapat dipercaya dalam
peluang.

\subsection*{Masalah Ulang Tahun}\index{masalah ulang tahun}

\begin{example}\label{exam 3.3} Berapa banyak orang yang perlu berada dalam sebuah
ruangan agar kita layak bertaruh bahwa dua orang di ruangan itu memiliki tanggal ulang
tahun yang sama (yakni, peluang keberhasilannya lebih besar daripada 1/2)?

Karena terdapat 365 tanggal ulang tahun yang mungkin, kita mungkin tergoda untuk
menebak bahwa diperlukan sekitar 1/2 dari jumlah tersebut, yaitu 183.  Dengan jumlah
itu Anda pasti memenangkan taruhan.  Sebenarnya, jumlah yang diperlukan agar taruhan
menguntungkan hanya 23.  Untuk menunjukkannya, kita mencari peluang $p_r$ bahwa dalam
ruangan berisi $r$ orang tidak ada tanggal ulang tahun yang sama; taruhan kita
menguntungkan jika peluang ini kurang dari setengah.

Anggaplah terdapat 365 tanggal ulang tahun yang mungkin bagi setiap orang (kita
mengabaikan tahun kabisat).  Urutkan orang-orang tersebut dari~1 hingga~$r$.  Untuk
sebuah titik sampel $\omega$, kita memilih suatu urutan ulang tahun yang mungkin dengan
panjang $r$, yang masing-masing dipilih dari salah satu di antara 365 tanggal yang
mungkin.  Terdapat 365 kemungkinan untuk unsur pertama dalam urutan, dan untuk setiap
pilihan tersebut terdapat 365 kemungkinan bagi unsur kedua, dan seterusnya, sehingga
menghasilkan $365^r$ urutan ulang tahun yang mungkin.  Kita harus mencari banyaknya
urutan yang tidak memuat tanggal ulang tahun yang sama.  Untuk urutan seperti itu, kita
dapat memilih sembarang dari 365 hari untuk unsur pertama, lalu sembarang dari 364 hari
yang tersisa untuk unsur kedua, 363 untuk unsur ketiga, dan seterusnya, sampai kita
membuat $r$ pilihan.  Untuk pilihan ke-$r$, akan terdapat $365 - r + 1$ kemungkinan.
Jadi, jumlah seluruh urutan tanpa tanggal ulang tahun yang sama adalah
$$ 365 \cdot 364 \cdot 363 \cdot \dots \cdot (365 - r + 1)\ .
$$ Dengan demikian, jika setiap urutan sama mungkinnya,
$$ p_r = \frac{365 \cdot 364 \cdot \dots \cdot (365 - r + 1)}{365^r}\ .
$$ Kita menyatakan hasil kali
$$(n)(n-1)\cdots (n - r +1)$$ dengan $(n)_r$ (dibaca ``$n$ turun $r$," atau ``$n$
lebih rendah $r$").  Jadi,
$$ p_r = \frac{(365)_r}{(365)^r}\ .
$$

Program {\bf Birthday}\index{Birthday (program)} melakukan perhitungan ini dan mencetak
peluang untuk $r = 20$ hingga~25.  Dengan menjalankan program tersebut, kita memperoleh
hasil yang ditunjukkan pada Tabel~\ref{table
3.1}.
\begin{table}
\centering
\begin{tabular}{cc} Jumlah orang     & Peluang semua tanggal ulang tahun
berbeda \\ 
\\
20 & .5885616 \\ 21 & .5563117 \\ 22 & .5243047 \\ 23 & .4927028 \\
24 & .4616557 \\ 25 & .4313003 \\ 
\end{tabular}
\caption{Masalah ulang tahun.}
\label{table 3.1}
\end{table}
Seperti yang kita nyatakan di atas, peluang bahwa tidak ada tanggal ulang tahun yang
sama berubah dari lebih besar daripada setengah menjadi kurang daripada setengah ketika
jumlah orang bertambah dari~22 menjadi~23.  Untuk melihat betapa kecil kemungkinan kita
kalah dalam taruhan pada jumlah orang yang lebih besar, kita menjalankan kembali program
tersebut dengan mencetak nilai dari $r = 10$ hingga $r = 100$ dalam langkah sebesar~10.
Kita melihat bahwa di ruangan berisi 40 orang, rasio peluang sudah sangat mendukung
terjadinya kesamaan tanggal ulang tahun, dan di ruangan berisi 100 orang, rasio peluang
tersebut sangat besar.
\begin{table}
\centering
\begin{tabular}{cc} Jumlah orang     & Peluang semua tanggal ulang tahun
berbeda \\ 
\\
10 & .8830518 \\ 20 & .5885616 \\ 30 & .2936838 \\ 40 & .1087682 \\
50 & .0296264 \\ 60 & .0058773 \\ 70 & .0008404 \\ 80 & .0000857 \\ 90 & .0000062 \\
100 & .0000003 \\ 
\end{tabular}
\caption{Masalah ulang tahun.}
\label{table 3.2}
\end{table}
Kita telah menganggap bahwa ulang tahun sama mungkinnya jatuh pada hari mana pun.
Bukti statistik menunjukkan bahwa anggapan ini tidak benar.  Namun, secara intuitif
jelas (meskipun tidak mudah dibuktikan) bahwa ketidakseragaman ini justru membuat
kesamaan ulang tahun lebih mungkin terjadi dalam kelompok berisi 23 orang.  (Lihat
Latihan~\ref{exer 3.1.19} untuk mengetahui apa yang terjadi di planet yang memiliki
lebih atau kurang dari 365 hari dalam setahun.)
\end{example}

Sekarang kita beralih ke topik permutasi.

\subsection*{Permutasi}\index{permutasi}
\begin{definition}\label{def 3.1} Misalkan $A$ suatu himpunan berhingga.  Sebuah
{\em permutasi dari $A$} adalah pemetaan satu-ke-satu dari $A$ ke dirinya sendiri.
\end{definition}

Untuk menentukan suatu permutasi tertentu, kita menuliskan unsur-unsur $A$ dan, di
bawahnya, menunjukkan ke mana setiap unsur dipetakan oleh pemetaan satu-ke-satu.
Sebagai contoh, jika $A = \{a,b,c\}$, salah satu permutasi $\sigma$ yang mungkin adalah
$$
\sigma = \pmatrix{ a & b & c \cr b & c & a \cr}.
$$

Oleh permutasi $\sigma$, $a$ dipetakan ke $b$, $b$ dipetakan ke $c$, dan $c$ dipetakan
ke $a$.  Syarat bahwa pemetaan itu satu-ke-satu berarti tidak ada dua unsur $A$ yang
dipetakan ke unsur $A$ yang sama.

Kita dapat menyusun unsur-unsur himpunan dalam suatu urutan dan menamainya kembali
1,~2, \dots,~$n$.  Kemudian, suatu permutasi khas dari himpunan $A =
\{a_1,a_2,a_3,a_4\}$ dapat ditulis dalam bentuk
$$
\sigma = \pmatrix{ 1 & 2 & 3 & 4 \cr 2 & 1 & 4 & 3 \cr},
$$ yang menunjukkan bahwa $a_1$ dipetakan ke $a_2$, $a_2$ ke $a_1$, $a_3$ ke $a_4$,
dan $a_4$ ke $a_3$.

Jika kita selalu memilih baris atas sebagai 1 2 3 4, untuk menentukan permutasi kita
hanya perlu memberikan baris bawah, dengan pengertian bahwa baris ini menunjukkan ke
mana 1, 2, dan seterusnya dipetakan.  Jika ditulis dengan cara ini, permutasi sering
disebut {\em penyusunan ulang} atas $n$ objek 1,~2,~3, \dots,~$n$.  Sebagai contoh,
semua permutasi, atau penyusunan ulang, yang mungkin atas himpunan $A = \{1,2,3\}$ adalah:
$$ 
123,\ 132,\ 213,\ 231,\ 312,\ 321\ .
$$
\par
Menghitung banyaknya permutasi yang mungkin atas $n$ objek merupakan hal yang mudah.
Menurut prinsip pencacahan umum kita, terdapat $n$ cara untuk menetapkan unsur pertama;
untuk setiap cara tersebut terdapat $n - 1$ cara untuk menetapkan objek kedua,
$n - 2$ untuk yang ketiga, dan seterusnya.  Hal ini membuktikan teorema berikut.

\begin{theorem}\label{thm 3.1}\label{thm 3.1.1} Jumlah seluruh permutasi dari suatu
himpunan $A$ yang memiliki $n$ unsur diberikan oleh $n
\cdot (n~-~1) \cdot (n - 2) \cdot \ldots \cdot 1$.
\end{theorem}

Kadang-kadang ada gunanya mempertimbangkan pengurutan subhimpunan dari suatu himpunan.
Hal ini mendorong definisi berikut.
\begin{definition} Misalkan $A$ suatu himpunan dengan $n$ unsur, dan misalkan $k$
suatu bilangan bulat antara 0 dan $n$.  Maka permutasi-$k$ dari $A$ adalah daftar
terurut suatu subhimpunan $A$ berukuran $k$.
\end{definition}

Dengan menggunakan teknik yang sama seperti dalam teorema sebelumnya, hasil berikut
mudah dibuktikan.

\begin{theorem}\label{thm 3.2} Jumlah seluruh permutasi-$k$ dari suatu himpunan $A$
yang memiliki $n$ unsur diberikan oleh
$n \cdot (n-1) \cdot (n-2) \cdot \ldots \cdot (n - k + 1)$.
\end{theorem}

\subsection*{Faktorial}

Bilangan yang diberikan dalam Teorema~\ref{thm 3.1} disebut {\em faktorial $n$,}
\index{faktorial} dan dinyatakan dengan $n!$\index{$n"!$}.  Ungkapan 0! didefinisikan
sebagai 1 agar rumus-rumus tertentu menjadi lebih sederhana.  Beberapa nilai pertama
fungsi ini ditunjukkan pada Tabel~\ref{table 3.25}.  Pembaca akan melihat bahwa fungsi
ini tumbuh sangat cepat.

\begin{table}
\centering
\begin{tabular}{rr} $n$     & $n!$ \\
\\ 
0 & 1 \\
1 & 1 \\
2 & 2 \\
3 & 6 \\
4 & 24 \\
5 & 120 \\
6 & 720 \\
7 & 5040 \\
8 & 40320 \\
9 & 362880 \\
10 & 3628800 \\

\end{tabular}
\caption{Nilai fungsi faktorial.}
\label{table 3.25}
\end{table}




\par Ungkapan $n!$ akan muncul dalam banyak perhitungan kita, dan kita akan memerlukan
perkiraan besarnya ketika $n$ besar.  Jelas tidak praktis melakukan perhitungan eksak
dalam hal ini.  Sebagai gantinya, kita akan menggunakan hasil yang disebut {\em rumus
Stirling.}  Sebelum menyatakan rumus tersebut, kita memerlukan sebuah definisi.

\begin{definition}\label{def 3.2} Misalkan $a_n$ dan $b_n$ dua barisan bilangan.  Kita
mengatakan bahwa $a_n$ {\em setara secara asimtotik\index{setara secara asimtotik}
dengan $b_n$}, dan menulis $a_n \sim b_n$, jika
$$
\lim_{n \to \infty} \frac{a_n}{b_n} = 1\ .
$$
\end{definition}
\begin{example}\label{exam 3.4} Jika $a_n = n + \sqrt n$ dan $b_n = n$, karena
$a_n/b_n = 1 + 1/\sqrt n$ dan rasio ini menuju 1 ketika $n$ menuju tak hingga, kita
memperoleh $a_n \sim b_n$.
\end{example}

\begin{theorem}{\bf (Rumus Stirling)}\label{thm 3.3}\index{rumus Stirling} Barisan
$n!$ setara secara asimtotik dengan
$$ n^ne^{-n}\sqrt{2\pi n}\ .
$$
\end{theorem}

Bukti rumus Stirling dapat ditemukan dalam sebagian besar buku analisis.  Mari kita
periksa pendekatan ini dengan menggunakan komputer.  Program {\bf
StirlingApproximations}\index{StirlingApproximations\\ (program)} mencetak
$n!$, pendekatan Stirling, dan terakhir rasio kedua bilangan tersebut.  Contoh keluaran
program ini ditunjukkan pada Tabel~\ref{table 3.26}.  Perhatikan bahwa meskipun rasio
kedua bilangan makin mendekati 1, selisih antara nilai eksak dan pendekatannya justru
meningkat.  Bahkan, selisih ini akan menuju tak hingga ketika $n$ menuju tak hingga,
meskipun rasionya menuju 1.  (Hal yang sama juga berlaku dalam Contoh~\ref{exam 3.4},
ketika $n +
\sqrt n \sim n$, tetapi selisihnya adalah $\sqrt n$.)

\begin{table}
\centering
\begin{tabular}{rrrr} $n$     & $n!$ &Pendekatan &Rasio \\
\\ 
1 & 1 & .922 & 1.084\\
2 & 2  & 1.919 & 1.042\\
3 & 6  & 5.836 & 1.028\\
4 & 24 & 23.506 & 1.021\\
5 & 120 & 118.019 & 1.016\\
6 & 720 & 710.078 & 1.013\\
7 & 5040 & 4980.396 & 1.011\\
8 & 40320 & 39902.395 & 1.010\\
9 & 362880 & 359536.873 & 1.009\\
10 & 3628800 & 3598696.619 & 1.008\\

\end{tabular}
\caption{Pendekatan Stirling untuk fungsi faktorial.}
\label{table 3.26}
\end{table}

\subsection*{Menghasilkan Permutasi Acak}
Sekarang kita membahas cara menghasilkan permutasi acak dari bilangan bulat antara 1 dan
$n$.  Perhatikan percobaan berikut.  Kita mulai dengan setumpuk $n$ kartu yang diberi
label 1 hingga $n$.  Kita memilih sebuah kartu secara acak dari tumpukan, mencatat
labelnya, lalu menyisihkan kartu itu.  Proses ini diulangi sampai seluruh $n$ kartu
terpilih.  Jelas bahwa setiap permutasi bilangan bulat dari 1 hingga $n$ dapat muncul
sebagai urutan label dalam percobaan ini, dan setiap urutan label sama mungkinnya untuk
muncul.  Dalam implementasi algoritma komputer kita, prosedur di atas disebut {\bf
RandomPermutation}.\index{RandomPermutation (program)}  



\subsection*{Titik Tetap} Ada banyak masalah menarik yang berkaitan dengan sifat
permutasi yang dipilih secara acak dari himpunan semua permutasi atas suatu himpunan
berhingga.  Sebagai contoh, karena permutasi merupakan pemetaan satu-ke-satu dari
himpunan ke dirinya sendiri, menarik untuk menanyakan berapa banyak titik yang dipetakan
ke dirinya sendiri.  Titik-titik seperti itu kita sebut {\em titik tetap}
\index{permutasi!titik tetap dari}\index{titik tetap} dari pemetaan tersebut.
\par
Misalkan $p_k(n)$ adalah peluang bahwa suatu permutasi acak dari himpunan
$\{1, 2, \ldots, n\}$ memiliki tepat $k$ titik tetap.  Kita akan mencoba mempelajari
peluang-peluang ini dengan menggunakan simulasi.  Program {\bf FixedPoints}
\index{FixedPoints (program)} menggunakan prosedur {\bf RandomPermutation} untuk
menghasilkan permutasi acak dan menghitung titik tetap.  Program ini mencetak proporsi
kejadian dengan $k$ titik tetap serta rata-rata jumlah titik tetap.  Hasil program untuk
500 simulasi pada kasus $n = 10$,~20, dan~30 ditunjukkan pada Tabel~\ref{table 3.3}.
\begin{table}
\centering
\begin{tabular}{|c|l|l|l|} \hline
Jumlah titik tetap & \multicolumn{3}{c|} {Proporsi permutasi} \\ \cline{2-4}
  & n = 10 & n = 20 & n = 30 \\ \hline
0 & \makebox[.375in][r]{.362}   & \makebox[.375in][r]{.370}    & \makebox[.375in][r]{.358}  
\\ 1 & \makebox[.375in][r]{.368}   & \makebox[.375in][r]{.396}   &
\makebox[.375in][r]{.358}  
\\ 2 & \makebox[.375in][r]{.202}   & \makebox[.375in][r]{.164}   &
\makebox[.375in][r]{.192}  
\\ 3 & \makebox[.375in][r]{.052}   & \makebox[.375in][r]{.060}    &
\makebox[.375in][r]{.070}   
\\ 4 & \makebox[.375in][r]{.012}   & \makebox[.375in][r]{.008}   &
\makebox[.375in][r]{.020}   
\\ 5 & \makebox[.375in][r]{.004}   & \makebox[.375in][r]{.002}   & \makebox[.375in][r]{.002}  
\\ \hline Rata-rata jumlah titik tetap & \makebox[.375in][r]{.996}   & \makebox[.375in][r]{.948}   
& \makebox[.125in][r]1.042 \\ \hline
\end{tabular}
\caption{Distribusi titik tetap.}
\label{table 3.3}
\end{table}
Perhatikan fakta yang cukup mengejutkan bahwa perkiraan peluang kita tampaknya tidak
terlalu bergantung pada banyaknya unsur dalam permutasi.  Sebagai contoh, peluang tidak
adanya titik tetap ketika $n = 10,\ 20,$ atau~30 diperkirakan berada antara
.35~dan~.37.  Kelak kita akan melihat (lihat Contoh~\ref{exam 3.13}) bahwa untuk
$n \geq 10$, peluang eksak $p_n(0)$, hingga ketelitian enam tempat desimal, sama dengan
$1/e \approx
.367879$.  Jadi, untuk semua keperluan praktis, setelah $n = 10$ peluang
bahwa suatu permutasi acak dari himpunan $\{1, 2, \ldots, n\}$ tidak memiliki titik
tetap tidak bergantung pada $n$.  Simulasi ini juga menyarankan bahwa rata-rata jumlah
titik tetap mendekati 1.  Dapat ditunjukkan (lihat Contoh~\ref{exam 6.1}) bahwa rata-rata
tersebut tepat sama dengan 1 untuk semua $n$.
\par
Versi masalah titik tetap yang lebih mudah dibayangkan antara lain: Anda telah menyusun
buku di rak menurut urutan abjad nama penulis, lalu buku-buku itu dikembalikan ke rak
secara acak; berapakah peluang bahwa tepat $k$ buku berakhir di posisi yang benar?
(Masalah perpustakaan.)\index{masalah perpustakaan}  Di sebuah restoran, $n$ topi
dititipkan lalu tercampur aduk; berapakah peluang tidak seorang pun
mendapatkan kembali topinya sendiri?  (Masalah penitipan topi.)
\index{masalah penitipan topi}  Dalam Catatan Historis di akhir bagian ini, kita
memberikan satu metode untuk menyelesaikan masalah penitipan topi secara eksak.  Metode
lain diberikan dalam Contoh~\ref{exam 3.13}.  


\subsection*{Rekor}\index{rekor}

Berikut adalah masalah peluang menarik lainnya yang melibatkan permutasi.  Perkiraan
jumlah curah salju yang diukur dalam inci di Hanover\index{salju di Hanover}, New
Hampshire, selama sepuluh tahun dari 1974 hingga 1983 ditunjukkan pada
Tabel~\ref{table 3.4}.
\begin{table}
\centering
\begin{tabular}{lc} Tahun & Curah salju dalam inci \\ \hline
  1974  & \hspace{.1in}75   \\
  1975  & \hspace{.1in}88   \\
  1976  & \hspace{.1in}72   \\
  1977  & \hspace{.031in}110 \\
  1978  & \hspace{.1in}85   \\
  1979  & \hspace{.1in}30   \\
  1980  & \hspace{.1in}55   \\
  1981  & \hspace{.1in}86   \\
  1982  & \hspace{.1in}51   \\
  1983  & \hspace{.1in}64   \\
\end{tabular}
\caption{Curah salju di Hanover.}
\label{table 3.4}
\end{table}
Misalkan kita mulai mencatat rekor pada 1974.  Curah salju pada tahun pertama dapat
dianggap sebagai rekor curah salju sejak pencatatan dimulai.  Rekor baru tercipta pada
1975; rekor berikutnya tercipta pada 1977, dan tidak ada rekor baru setelah tahun itu.
Jadi, dalam periode sepuluh tahun ini tercipta tiga rekor: 1974, 1975, dan 1977.
Pertanyaan kita adalah: Berapa banyak rekor yang dapat kita harapkan tercipta dalam
periode sepuluh tahun seperti itu?  Kita dapat menghitung banyaknya rekor dalam bentuk
suatu permutasi sebagai berikut: Kita menomori tahun dari 1 hingga~10.  Besarnya curah
salju yang sebenarnya tidak penting, tetapi ukuran relatifnya penting.  Karena itu,
kita dapat mengubah angka curah salju menjadi bilangan 1 hingga~10 dengan mengganti
angka terkecil dengan 1, angka terkecil berikutnya dengan 2, dan seterusnya.  (Kita
menganggap tidak ada nilai yang sama.)  Untuk contoh ini, kita memperoleh data yang
ditunjukkan pada Tabel~\ref{table 3.5}.

\begin{table}
\centering
\begin{tabbing}
\hskip.5in  Tahun \hskip.5in \= 1\hskip.3in\= 2\hskip.3in \= 3\hskip.3in \= \,\,4\hskip.3in 
                  \= 5\hskip.3in\= 6\hskip.3in \= 7\hskip.3in \= 8\hskip.3in \=
9\hskip.3in \= 10\\
\hskip.5in  Peringkat        \> 6 \> 9 \> 5 \> 10 \> 7 \> 1 \> 3 \> 8 \> 2 \> \,\,4 
\end{tabbing}
\caption{Peringkat curah salju total.}
\label{table 3.5}
\end{table}

Ini memberi kita suatu permutasi bilangan dari 1~hingga~10, dan dari permutasi ini kita
dapat membaca rekor-rekornya; rekor tersebut terjadi pada tahun 1,~2, dan~4.  Dengan
demikian, kita dapat mendefinisikan rekor untuk suatu permutasi sebagai berikut:

\begin{definition}\label{def 3.3} Misalkan $\sigma$ suatu permutasi dari himpunan
$\{1, 2, \ldots, n\}$.  Maka $i$ adalah suatu {\em rekor} dari $\sigma$ jika
$i = 1$ atau $\sigma(j) < \sigma(i)$ untuk setiap $j = 1,\ldots,\,i - 1$.
\end{definition}

Sekarang, jika kita menganggap semua peringkat curah salju selama periode $n$ tahun
sama mungkinnya (dan tidak mengizinkan nilai yang sama), kita dapat memperkirakan
peluang terdapat $k$ rekor dalam $n$ tahun serta rata-rata jumlah rekor melalui simulasi.
\par
Kita telah menulis program {\bf Records}\index{Records (program)} yang menghitung jumlah
rekor dalam permutasi yang dipilih secara acak.  Kita telah menjalankan program ini
untuk kasus $n = 10$,~20,~30.  Untuk $n = 10$, rata-rata jumlah rekor adalah 2.968,
untuk 20 nilainya 3.656, dan untuk 30 nilainya 3.960.  Sekarang terlihat bahwa rata-rata
tersebut meningkat, tetapi sangat lambat.  Kelak kita akan melihat (lihat
Contoh~\ref{exam 6.5}) bahwa rata-rata itu kira-kira $\log n$.  Karena
$\log 10 = 2.3$,
$\log 20 = 3$, dan
$\log 30 = 3.4$, hal ini konsisten dengan hasil simulasi kita.
\par
Seperti telah disebutkan, kita akan dapat memperoleh rumus untuk hasil eksak beberapa
masalah sejenis di atas.  Namun, perubahan kecil saja pada masalah dapat membuat hal
ini mustahil.  Kekuatan simulasi adalah bahwa perubahan kecil pada suatu masalah tidak
membuat simulasinya jauh lebih sulit.  (Lihat Latihan~\ref{exer 3.1.20} untuk variasi
menarik dari masalah penitipan topi.)

\subsection*{Daftar Permutasi}

Metode lain untuk menyelesaikan masalah yang tidak peka terhadap perubahan kecil adalah
meminta komputer mencantumkan semua permutasi yang mungkin dan menghitung proporsi yang
memiliki sifat yang diinginkan.  Program {\bf AllPermutations}
\index{AllPermutations (program)} menghasilkan daftar seluruh permutasi dari $n$.
Ketika mencoba menjalankan program ini, kita berhadapan dengan keterbatasan penggunaan
komputer.  Banyaknya permutasi dari $n$ meningkat sedemikian cepat sehingga bahkan
mencantumkan semua permutasi dari 20 objek pun tidak praktis.

\subsection*{Catatan Historis} 
Prinsip pencacahan dasar kita menyatakan bahwa jika Anda dapat melakukan satu hal
dengan $r$ cara dan, untuk setiap cara tersebut, melakukan hal lain dengan $s$ cara,
maka Anda dapat melakukan kedua hal tersebut dengan $rs$ cara.  Hasil ini begitu jelas
dengan sendirinya sehingga Anda mungkin menduga hasil tersebut muncul sangat awal dalam
matematika.  N.~L. Biggs menyarankan agar kita menelusuri salah satu contoh prinsip ini
sebagai berikut: Pertama, ia mengisahkan sebuah sajak anak-anak\index{sajak anak-anak}
populer yang berasal setidaknya dari 1730:\index{St. Ives}
\par
\vskip .2in
\begin{tabular}{@{\extracolsep{1in}}ll}
&Ketika aku pergi ke St. Ives,\\
&Aku bertemu pria dengan tujuh istri,\\
&Setiap istri memiliki tujuh karung,\\
&Setiap karung memiliki tujuh kucing,\\
&Setiap kucing memiliki tujuh anak kucing.\\
&Anak kucing, kucing, karung, dan istri,\\
&Berapa banyak yang pergi ke St. Ives?
\end{tabular}
\par
\vskip .2in
\noindent
(Anda hanya memerlukan prinsip kita jika tidak cukup cermat untuk menyadari bahwa
jawabannya seharusnya {\em satu,} sebab hanya penuturnya yang sedang pergi ke
St.~Ives; yang lain berjalan ke arah berlawanan!)
\par Ia juga memberikan suatu masalah yang muncul dalam salah satu naskah matematika
tertua yang masih bertahan, dari sekitar 1650~{\footnotesize B.C.}, yang secara kasar
diterjemahkan sebagai berikut:
\par
\vskip .2in
\begin{tabular}{lr}
 \hspace{1.5in}Rumah \hspace{2in}&   7\\
 \hspace{1.5in}Kucing           &  49                  \\
 \hspace{1.5in}Tikus            &  343                 \\
 \hspace{1.5in}Gandum           &  2401         \\
 \hspace{1.5in}Hekat            &  {\underbar{16807}} \\
                  &  19607              \\
\end{tabular}
\par
\vskip .2in
Penafsiran berikut telah diajukan: terdapat tujuh rumah, masing-masing dengan tujuh
kucing; setiap kucing membunuh tujuh tikus; setiap tikus akan memakan tujuh bulir
gandum, yang masing-masing akan menghasilkan tujuh takaran hekat biji-bijian.  Dengan
penafsiran ini, tabel tersebut menjawab pertanyaan tentang berapa banyak takaran hekat
yang diselamatkan oleh tindakan kucing-kucing itu.  Tidak jelas mengapa penulis tabel
ingin menjumlahkan bilangan-bilangan tersebut.\index{BIGGS, N. L.}\footnote{N.~L. Biggs, ``The Roots of Combinatorics," {\em
Historia Mathematica,} vol.~6 (1979), pp.~109--136.}

Salah satu penggunaan faktorial paling awal muncul dalam bukti Euclid\index{EUCLID}
bahwa terdapat tak berhingga banyaknya bilangan prima.  Euclid berpendapat bahwa pasti
ada bilangan prima di antara $n$ dan $n! + 1$ sebagai berikut: $n!$ dan $n! + 1$ tidak
dapat memiliki faktor persekutuan.  Bilangan $n! + 1$ merupakan bilangan prima atau
memiliki faktor sejati.  Dalam kasus kedua, faktor tersebut tidak dapat membagi $n!$
dan karenanya harus berada di antara $n$ dan $n! + 1$.  Jika faktor ini bukan bilangan
prima, faktor tersebut memiliki faktor yang, dengan argumen yang sama, harus lebih besar
daripada $n$.  Dengan cara ini, pada akhirnya kita mencapai bilangan prima yang lebih
besar daripada $n$, dan hal ini berlaku untuk setiap $n$.

Kaidah ``$n!$" untuk banyaknya permutasi tampaknya pertama kali muncul di India.
Contoh-contohnya telah ditemukan sejak 300~{\footnotesize B.C.}, dan pada abad kesebelas
rumus umumnya tampaknya telah dikenal luas di India, kemudian di negeri-negeri Arab.

{\em Masalah penitipan topi}\index{masalah penitipan topi} ditemukan dalam sebuah buku
peluang awal yang ditulis oleh de~Montmort\index{de MONTMORT, P. R.} dan pertama kali
dicetak pada 1708.\footnote{P.~R. de~Montmort, {\em
Essay d'Analyse sur des Jeux de Hazard,} 2d~ed. (Paris: Quillau, 1713).}  Buku tersebut
memuat permainan bernama {\em Treize.}\index{Treize}  Dalam versi sederhana permainan ini yang
dibahas oleh de~Montmort, seseorang membalik kartu bernomor 1 hingga 13 sambil menyebut
1,~2, \dots,~13 ketika kartu-kartu diperiksa.  De~Montmort menanyakan peluang bahwa
tidak ada kartu yang dibalik cocok dengan nomor yang disebutkan.
\par
Peluang ini sama dengan peluang bahwa suatu permutasi acak atas 13 unsur tidak memiliki
titik tetap.  De~Montmort menyelesaikan masalah ini dengan menggunakan relasi rekursi
sebagai berikut: misalkan $w_n$ adalah banyaknya permutasi atas $n$ unsur tanpa titik
tetap (permutasi demikian disebut {\em permutasi tanpa titik tetap})
\index{permutasi tanpa titik tetap}.  Maka $w_1 = 0$ dan
$w_2 = 1$.
\par Sekarang anggaplah $n \ge 3$ dan pilih suatu permutasi tanpa titik tetap atas
bilangan bulat antara 1 dan $n$.  Misalkan $k$ adalah bilangan bulat pada posisi pertama
dalam permutasi tersebut.  Menurut definisi permutasi tanpa titik tetap, kita memiliki
$k \ne 1$.  Ada dua kemungkinan penting mengenai posisi 1 dalam permutasi: 1 berada
pada posisi ke-$k$ atau berada di tempat lain.  Dalam kasus pertama, $n-2$ bilangan bulat
yang tersisa dapat ditempatkan dengan $w_{n-2}$ cara tanpa menghasilkan titik tetap.
Dalam kasus kedua, kita mempertimbangkan himpunan bilangan bulat
$\{1, 2, \ldots, k-1, k+1, \ldots, n\}$.  Bilangan-bilangan dalam himpunan ini harus
menempati posisi $\{2, 3, \ldots, n\}$ sedemikian rupa sehingga tidak satu pun bilangan
selain 1 dalam himpunan tersebut menjadi titik tetap, dan juga agar 1 tidak berada pada
posisi $k$.  Banyaknya cara untuk menghasilkan susunan seperti ini tepat sebesar
$w_{n-1}$.  Karena terdapat $n-1$ nilai $k$ yang mungkin, kita memperoleh
$$ w_n = (n - 1)w_{n - 1} + (n - 1)w_{n -2}
$$ untuk $n \ge 3$.  Dari persamaan terakhir ini, kita mungkin menduga bahwa barisan
$\{w_n\}$ tumbuh seperti barisan $\{n!\}$.  
\par Sebenarnya, dengan induksi mudah dibuktikan bahwa
$$ w_n = nw_{n - 1} + (-1)^n\ .
$$ Maka $p_i = w_i/i!$ memenuhi
$$ p_i - p_{i - 1} = \frac{(-1)^i}{i!}\ .
$$ Jika kita menjumlahkan dari $i = 2$ hingga~$n$, dan menggunakan fakta bahwa
$p_1 = 0$, kita memperoleh
$$ p_n = \frac1{2!} - \frac1{3!} + \cdots + \frac{(-1)^n}{n!}\ .
$$ Ini sama dengan $n + 1$ suku pertama dalam ekspansi $e^x$ untuk $x = -1$, sehingga
untuk $n$ besar nilainya kira-kira $e^{-1} \approx .368$.  David
\index{DAVID, F. N.} menyatakan bahwa ini mungkin merupakan penggunaan pertama fungsi
eksponensial dalam peluang.\footnote{F.~N. David, {\em Games, Gods and Gambling} (London: Griffin,
1962), p.~146.}  Pada bagian berikutnya, kita akan melihat cara lain untuk menurunkan
hasil de~Montmort dengan menggunakan metode yang dikenal sebagai metode inklusi-eksklusi.
\par
Baru-baru ini, sebuah masalah terkait muncul dalam kolom Marilyn vos Savant.\footnote{M.
vos Savant, Ask Marilyn, \emx{Parade Magazine, Boston Globe}, 21
August 1994.}\index{vos SAVANT, M.}  Charles Price\index{PRICE, C.} menulis untuk
menanyakan pengalamannya memainkan suatu bentuk solitaire yang kadang-kadang disebut
``frustration solitaire."\index{frustration solitaire}  Dalam permainan ini, setumpuk
kartu dikocok, lalu dibagikan satu per satu.  Ketika kartu dibagikan, pemain menghitung
dari 1 hingga 13, kemudian mulai lagi dari 1.  (Dengan demikian, setiap bilangan disebut
empat kali.)  Jika bilangan yang sedang disebut sama dengan nilai kartu yang sedang
dibuka, pemain kalah.  Price mendapati bahwa ia jarang menang dan bertanya-tanya seberapa
sering semestinya ia menang.  Vos~Savant menyatakan bahwa jumlah kecocokan yang diharapkan
adalah 4, sehingga permainan tersebut seharusnya sulit dimenangkan.
\par
Mencari peluang menang merupakan masalah yang lebih sulit daripada masalah yang
diselesaikan de~Montmort karena ketika seluruh tumpukan kartu diperiksa, terdapat pola
kecocokan berbeda yang mungkin terjadi.  Sebagai contoh, kecocokan dapat terjadi pada
dua kartu dengan nilai sama, misalnya dua kartu as, atau pada dua nilai berbeda,
misalnya kartu dua dan kartu tiga.
\par
Pembahasan masalah ini dapat ditemukan dalam Riordan.\footnote{J. Riordan, \emx{ An
Introduction to Combinatorial Analysis,} (New York: John Wiley \& Sons, 1958).}\index{RIORDAN, J.}
Dalam buku tersebut ditunjukkan bahwa ketika $n \rightarrow \infty$, peluang tidak
adanya kecocokan menuju $1/e^4$.
\par
Permainan Treize asli lebih sulit dianalisis daripada frustration solitaire.  Treize
dimainkan sebagai berikut.  Satu orang dipilih sebagai bandar dan yang lain menjadi
pemain.  Setiap pemain selain bandar memasang taruhan.  Bandar mengocok kartu dan
membukanya satu per satu sambil menyebut, ``As, dua, tiga,..., raja," seperti dalam
frustration solitaire.  Jika bandar melewati 13 kartu tanpa kecocokan, ia membayar para
pemain sejumlah taruhan mereka, lalu peran bandar berpindah kepada orang lain.  Jika
terjadi kecocokan, bandar mengambil taruhan para pemain; para pemain memasang taruhan
baru, dan bandar melanjutkan membuka kartu sambil menyebut, ``As, dua, tiga, ...."
Jika kehabisan kartu, bandar mengocok ulang dan melanjutkan hitungan dari tempat ia
berhenti.  Ia terus bermain sampai terdapat rangkaian 13 kartu tanpa kecocokan, lalu
bandar baru dipilih.
\par
Pertanyaannya sekarang adalah berapa banyak uang yang dapat diharapkan dimenangkan
bandar dari setiap pemain.  De~Montmort menemukan bahwa jika setiap pemain memasang
taruhan sebesar 1, misalnya, bandar akan menang kira-kira .801 dari setiap pemain.  
\par
Peter Doyle\index{DOYLE, P. G.} menghitung secara eksak besarnya kemenangan yang dapat
diharapkan oleh bandar.  Jawabannya adalah:
\vskip .2in
\noindent
26516072156010218582227607912734182784642120482136091446715371962089931\break
52311343541724554334912870541440299239251607694113500080775917818512013\break    
82176876653563173852874555859367254632009477403727395572807459384342747\break    
87664965076063990538261189388143513547366316017004945507201764278828306\break    
60117107953633142734382477922709835281753299035988581413688367655833113\break
24476153310720627474169719301806649152698704084383914217907906954976036\break
28528211590140316202120601549126920880824913325553882692055427830810368\break
57818861208758248800680978640438118582834877542560955550662878927123048\break
26997601700116233592793308297533642193505074540268925683193887821301442\break
70519791882/\break
33036929133582592220117220713156071114975101149831063364072138969878007\break
99647204708825303387525892236581323015628005621143427290625658974433971\break
65719454122908007086289841306087561302818991167357863623756067184986491\break
35353553622197448890223267101158801016285931351979294387223277033396967\break
79797069933475802423676949873661605184031477561560393380257070970711959\break
69641268242455013319879747054693517809383750593488858698672364846950539\break
88868628582609905586271001318150621134407056983214740221851567706672080\break
94586589378459432799868706334161812988630496327287254818458879353024498\break
00322425586446741048147720934108061350613503856973048971213063937040515\break
59533731591.\break
\par
\noindent
Nilai ini adalah .803 jika dibulatkan hingga 3 tempat desimal.  Uraian algoritma yang
digunakan untuk menemukan jawaban ini dapat ditemukan di halaman Web miliknya.\footnote{P.~Doyle, ``Solution to Montmort's Probleme du Treize,''
http://math.ucsd.edu/$\tilde{\ }$doyle/.}  Pembahasan masalah ini dan masalah lainnya
dapat ditemukan dalam Doyle et al.\footnote{P. Doyle, C. Grinstead, and J. Snell, ``Frustration Solitaire,"
\emx{UMAP Journal}, vol.\ 16, no.\ 2 (1995), pp.~137-145.}\index{GRINSTEAD, C. M.}\index{SNELL, J. L.}
\par
{\em Masalah ulang tahun} tampaknya tidak memiliki sejarah yang sangat panjang.  Masalah
jenis ini pertama kali dibahas oleh von~Mises.\index{von MISES, R.}\footnote{R. von~Mises, ``\"Uber
Aufteilungs- und Besetzungs-Wahrscheinlichkeiten," {\em Revue de la Facult\'e des
Sciences de l'Universit\'e d'Istanbul, N. S.} vol.~4 (1938-39), pp.~145-163.}  Masalah ini
dipopulerkan pada dekade 1950-an oleh buku Feller.\footnote{W.~Feller, \emx {Introduction to
Probability Theory and Its Applications,} vol.~1, 3rd~ed. (New York: John Wiley \&
Sons, 1968).}
\par
Stirling menyajikan rumusnya
$$ n! \sim \sqrt{2\pi n}\left(\frac{n}{e}\right)^n
$$ dalam karyanya {\em Methodus Differentialis} yang diterbitkan pada
1730.\index{STIRLING, J.}\footnote{J.~Stirling, {\em Methodus Differentialis,} (London: Bowyer,
1730).}  Pendekatan ini digunakan oleh de~Moivre ketika merumuskan teorema limit pusatnya
yang terkenal, yang akan kita pelajari dalam Bab~\ref{chp 9}.  De~Moivre sendiri telah
menetapkan pendekatan ini secara mandiri, tetapi tanpa mengidentifikasi konstanta $\pi$.
Setelah menetapkan pendekatan
$$
\frac{2B}{\sqrt n}
$$ untuk suku tengah distribusi binomial, dengan konstanta $B$ yang ditentukan oleh
suatu deret tak hingga, de~Moivre menulis:
\begin{quote}
\dots~sahabat saya yang terhormat dan terpelajar, Tn.\ James Stirling, yang setelah saya
menekuni penyelidikan itu, menemukan bahwa Besaran $B$ menyatakan Akar Kuadrat Keliling
Lingkaran yang Jari-jarinya Satu, sehingga jika Keliling itu disebut $c$, Rasio Suku
tengah terhadap Jumlah semua Suku akan dinyatakan oleh
$2/\sqrt{nc}\,$\dots.\index{de MOIVRE, A.}\footnote{A. de~Moivre, {\em The Doctrine of Chances,}
3rd~ed. (London: Millar, 1756).}
\end{quote}

\exercises
\begin{LJSItem}

\i\label{exer 3.1.1} Empat orang akan disusun dalam satu baris untuk difoto.  Berapa
banyak cara penyusunan yang dapat dilakukan?

\i\label{exer 3.1.2} Sebuah produsen mobil menyediakan empat warna untuk bagian luar
mobil dan tiga warna untuk bagian dalam.  Berapa banyak kombinasi warna berbeda yang
dapat diproduksinya?

\i\label{exer 3.1.3} Dalam komputer digital, sebuah {\em bit} adalah salah satu bilangan
bulat \{0,1\}, dan sebuah {\em kata} adalah sembarang untaian 32 bit.  Berapa banyak
kata berbeda yang mungkin?

\i\label{exer 3.1.4} Berapakah peluang bahwa setidaknya 2 presiden Amerika Serikat
meninggal pada tanggal yang sama dalam setahun?  Jika Anda bertaruh bahwa hal ini telah
terjadi, apakah Anda akan memenangkan taruhan?

\i\label{exer 3.1.5} Terdapat tiga rute berbeda yang menghubungkan kota A dengan kota B.
Berapa banyak cara yang dapat digunakan untuk melakukan perjalanan pulang-pergi dari A
ke B lalu kembali?
Berapa banyak cara jika rute pulang harus berbeda?

\i\label{exer 3.1.6} Dalam menyusun orang mengelilingi meja bundar, kita memperhitungkan
posisi tempat duduk mereka relatif satu sama lain, bukan posisi sebenarnya dari salah
satu orang.  Tunjukkan bahwa $n$ orang dapat disusun mengelilingi meja bundar dengan
$(n - 1)!$ cara.

\i\label{exer 3.1.7} Lima orang menaiki sebuah lift\index{lift} yang berhenti di lima
lantai.  Dengan menganggap setiap orang memiliki peluang yang sama untuk menuju lantai
mana pun, carilah peluang bahwa mereka semua turun di lantai yang berbeda.

\i\label{exer 3.1.8} Suatu himpunan berhingga $\Omega$ memiliki $n$ unsur.  Tunjukkan
bahwa jika himpunan kosong dan $\Omega$ dihitung sebagai subhimpunan, terdapat $2^n$
subhimpunan dari $\Omega$.

\i\label{exer 3.1.9} Ketaksamaan yang lebih teliti untuk mendekati $n!$ diberikan oleh
$$
\sqrt{2\pi n}\left(\frac ne\right)^n e^{1/(12n + 1)} < n! < \sqrt{2\pi n}\left(\frac
ne\right)^n e^{1/(12n)}\ .
$$ Tulislah program komputer untuk mengilustrasikan ketaksamaan ini bagi $n = 1$ hingga~9.

\i\label{exer 3.1.10} Setumpuk kartu biasa dikocok dan 13 kartu dibagikan.  Berapakah
peluang bahwa kartu terakhir yang dibagikan adalah kartu as?

\i\label{exer 3.1.11} Terdapat $n$ pelamar untuk jabatan direktur komputasi.  Setiap
anggota komite pencarian yang beranggotakan tiga orang mewawancarai para pelamar secara
mandiri dan memberi peringkat dari 1 hingga $n$.  Seorang kandidat akan dipekerjakan
jika setidaknya dua dari ketiga pewawancara menempatkannya pada peringkat pertama.
Carilah peluang bahwa seorang kandidat akan diterima jika para anggota komite sama
sekali tidak mampu menilai kandidat dan hanya memberi peringkat secara acak.  Secara
khusus, bandingkan peluang ini untuk kasus tiga kandidat dan kasus sepuluh kandidat.

\i\label{exer 3.1.12} Sebuah orkestra simfoni memiliki repertoar yang terdiri atas 30
simfoni Haydn, 15 karya modern, dan 9 simfoni Beethoven.  Programnya selalu terdiri
atas sebuah simfoni Haydn yang diikuti karya modern, lalu sebuah simfoni Beethoven.

\begin{enumerate}
\item Berapa banyak program berbeda yang dapat dimainkannya?

\item Berapa banyak program berbeda yang ada jika ketiga karya dapat dimainkan dalam
urutan apa pun?

\item Berapa banyak program tiga karya yang berbeda jika lebih dari satu karya dari
kategori yang sama dapat dimainkan dan karya-karya itu dapat dimainkan dalam urutan
apa pun?
\end{enumerate}

\i\label{exer 3.1.13} Suatu negara bagian memiliki pelat nomor kendaraan yang memuat
tiga angka dan tiga huruf.  Berapa banyak pelat nomor berbeda yang mungkin

\begin{enumerate}
\item jika angka harus mendahului huruf?

\item jika tidak ada batasan mengenai letak huruf dan angka?
\end{enumerate}

\i\label{exer 3.1.14} Pintu pusat komputer memiliki kunci dengan lima tombol yang
diberi nomor dari~1 hingga~5.  Kombinasi angka yang membuka kunci merupakan urutan
lima angka dan diatur ulang setiap minggu.

\begin{enumerate}
\item Berapa banyak kombinasi yang mungkin jika setiap tombol harus digunakan sekali?

\item Anggaplah kunci juga dapat memiliki kombinasi yang mengharuskan Anda menekan dua
tombol secara bersamaan, lalu tiga tombol lainnya satu per satu.  Berapa banyak kombinasi
tambahan yang dimungkinkan?
\end{enumerate}

\i\label{exer 3.1.15} Sebuah pusat komputasi memiliki 3 prosesor yang menerima $n$
pekerjaan.  Pekerjaan tersebut dialokasikan kepada prosesor sepenuhnya secara acak sehingga
semua $3^n$ alokasi yang mungkin sama mungkinnya.  Carilah peluang bahwa tepat satu
prosesor tidak menerima pekerjaan.

\i\label{exer 3.1.16} Buktikan bahwa setidaknya dua orang di Atlanta, Georgia, memiliki
inisial yang sama, dengan menganggap tidak seorang pun memiliki lebih dari empat inisial.

\i\label{exer 3.1.17} Carilah rumus untuk peluang bahwa di antara suatu himpunan berisi
$n$ orang, setidaknya dua orang berulang tahun pada bulan yang sama (dengan menganggap
ulang tahun sama mungkinnya jatuh pada setiap bulan).

\i\label{exer 3.1.18} Perhatikan masalah mencari peluang adanya lebih dari satu kesamaan
ulang tahun dalam kelompok berisi $n$ orang.  Hal ini mencakup, misalnya, tiga orang
dengan ulang tahun yang sama, atau dua pasang orang dengan ulang tahun yang sama, atau
kesamaan yang lebih besar.  Tunjukkan bagaimana Anda dapat menghitung peluang ini, lalu
tulislah program komputer untuk melakukan perhitungan tersebut.  Gunakan program Anda
untuk mencari jumlah orang terkecil sehingga bertaruh bahwa akan terjadi lebih dari satu
kesamaan ulang tahun merupakan taruhan yang menguntungkan.

\istar\label{exer 3.1.19} Misalkan di planet Zorg\index{Zorg, planet}, satu tahun
memiliki $n$ hari, dan makhluk hidup di sana sama mungkinnya menetas pada hari mana pun
dalam setahun.  Kita ingin memperkirakan $d$, yaitu jumlah minimum makhluk hidup yang
diperlukan agar peluang setidaknya dua di antaranya memiliki tanggal ulang tahun yang
sama melebihi 1/2.
\begin{enumerate}
\item Dalam Contoh~\ref{exam 3.3}, telah ditunjukkan bahwa dalam suatu himpunan berisi
$d$ makhluk hidup, peluang bahwa tidak ada dua makhluk hidup yang memiliki tanggal
ulang tahun yang sama
adalah $${{(n)_d}\over{n^d}}\ ,$$ dengan $(n)_d = (n)(n-1)\cdots (n-d+1)$. Jadi, kita
ingin menyamakannya dengan 1/2 dan menyelesaikan persamaan untuk $d$.

\item  Dengan menggunakan rumus Stirling, tunjukkan bahwa
$${{(n)_d}\over{n^d}} \sim \biggl(1 + {d\over{n-d}}\biggr)^{n-d + 1/2} e^{-d}\ .$$
\item Sekarang ambil logaritma dari ungkapan di ruas kanan, dan gunakan fakta bahwa
untuk nilai $x$ yang kecil, kita memiliki
$$\log(1+x) \sim x - {{x^2}\over 2}\ .$$ (Secara tersirat, kita menggunakan fakta bahwa
$d$ berorde lebih kecil daripada $n$.  Kita juga akan menggunakan fakta ini dalam bagian (d).)

\item Samakan ungkapan yang ditemukan dalam bagian (c) dengan $-\log(2)$, lalu
selesaikan untuk $d$ sebagai fungsi dari $n$, sehingga menunjukkan bahwa
$$d \sim \sqrt{2(\log 2)\,n}\ .$$ 
\emx {Petunjuk}:  Jika ketiga suku dalam ungkapan yang ditemukan pada bagian (b)
digunakan, diperoleh persamaan kubik dalam $d$.  Jika suku terkecil dari ketiganya
dibuang, diperoleh persamaan kuadrat dalam $d$.

\item Gunakan komputer untuk menghitung nilai eksak $d$ bagi berbagai nilai $n$.
Bandingkan nilai-nilai ini dengan nilai pendekatan yang diperoleh menggunakan jawaban
bagian (d).
\end{enumerate}

\i\label{exer 3.1.20}  Dalam sebuah konferensi matematika, sepuluh peserta duduk secara
acak mengelilingi meja bundar saat makan.  Dengan menggunakan simulasi, perkirakan
peluang bahwa tidak ada sepasang orang yang duduk bersebelahan saat makan siang sekaligus
saat makan malam.  Dapatkah Anda membuat dugaan yang masuk akal untuk kasus $n$ peserta
ketika $n$ besar?

\i\label{exer 3.1.21} Ubahlah program {\bf AllPermutations} untuk menghitung banyaknya
permutasi dari $n$ objek yang memiliki tepat $j$ titik tetap untuk $j = 0$,~1,~2,
\dots,~$n$.  Jalankan program Anda untuk $n = 2$ hingga~6.  Buatlah dugaan mengenai
hubungan antara banyaknya permutasi yang memiliki 0 titik tetap dan yang memiliki tepat
1 titik tetap.  Bukti dugaan yang benar dapat ditemukan dalam Wilf.
\index{WILF, H. S.}\footnote{H.~S.
Wilf, ``A Bijection in the Theory of Derangements," \emx {Mathematics Magazine,} vol.~57, no.~1
(1984), pp.~37--40. }

\i\label{exer 3.1.22} Tn.~Wimply Dimple, salah seorang pembuat jam paling bergengsi di
London, datang menemui Sherlock Holmes\index{Holmes, Sherlock} dalam keadaan panik setelah
mengetahui bahwa seseorang telah memproduksi dan menjual tiruan kasar dari jam tangan
terlarisnya\index{jam tangan, palsu}.  Sebanyak 16 barang palsu yang ditemukan sejauh ini
memiliki nomor cap yang semuanya berada di antara 1 dan 56, dan Dimple ingin sekali
mengetahui seberapa besar skala pekerjaan si pemalsu.  Semua yang hadir sepakat bahwa masuk
akal untuk menganggap barang palsu yang diproduksi sejauh ini memiliki nomor berurutan
dari~1 hingga berapa pun jumlah totalnya.
\par
``Tenanglah, Dimple," ujar Dr.\ Watson.  ``Jika saya menjadi Anda, saya tidak akan
terlalu khawatir; Prinsip Kemungkinan Maksimum\index{Prinsip Kemungkinan\\ Maksimum}, yang
memperkirakan jumlah total sebagai nilai yang memberikan peluang tertinggi bagi deret
nomor yang ditemukan, menyarankan agar kita menebak 56 itu sendiri sebagai jumlah total.
Jadi, para pemalsumu bukanlah jaringan besar, dan kita akan menjebloskan mereka ke balik
jeruji sebelum usahamu mengalami kerugian berarti."
\par
``Omong kosong; tak usah bawa-bawa prinsipmu yang muluk-muluk itu, Watson," balas
Holmes.  ``Siapa pun dapat melihat bahwa jumlah jam tangan itu tentu jauh lebih banyak
daripada 56---rasio peluang bahwa kita kebetulan menemukan tepat jam tangan dengan
nomor tertinggi yang diproduksi begitu kecil hingga menggelikan.  Tebakan yang jauh
lebih baik adalah
\emx{dua kali} 56."

\begin{enumerate}
\item Tunjukkan bahwa Watson benar: Prinsip Kemungkinan Maksimum memberikan 56.

\item Tulislah program komputer untuk membandingkan strategi tebakan Holmes dan Watson
sebagai berikut: tetapkan jumlah total $N$ dan pilih secara acak 16 bilangan bulat antara
1 dan $N$.  Misalkan $m$ menyatakan bilangan terbesar di antaranya.  Maka tebakan Watson
untuk $N$ adalah $m$, sedangkan tebakan Holmes adalah $2m$.  Lihat tebakan mana yang lebih
dekat dengan $N$.  Ulangi percobaan ini (dengan $N$ tetap) seratus kali atau lebih, lalu
tentukan proporsi percobaan ketika setiap tebakan lebih dekat.  Strategi siapa yang tampaknya
lebih baik?
\end{enumerate}

\i\label{exer 3.1.23} Barbara Smith sedang mewawancarai kandidat untuk menjadi
sekretarisnya.  Saat mewawancarai para kandidat, ia dapat menentukan peringkat relatif
mereka, tetapi bukan peringkat sebenarnya.  Jadi, jika terdapat enam kandidat dengan
peringkat sebenarnya 6, 1, 4, 2, 3, 5 (dengan 1 sebagai yang terbaik), setelah
mewawancarai tiga kandidat pertama ia akan memberi mereka peringkat 3, 1, 2.  Ketika
mewawancarai setiap kandidat, ia harus menerima atau menolaknya.  Jika tidak menerima
kandidat setelah wawancara, ia kehilangan kesempatan untuk memilih kandidat tersebut.
Ia ingin menentukan strategi untuk memutuskan kapan harus berhenti dan menerima seorang
kandidat agar peluang memperoleh kandidat terbaik menjadi maksimum.  Anggaplah terdapat
$n$ kandidat dan mereka tiba dalam urutan peringkat yang acak.

\begin{enumerate}
\item Berapakah peluang Barbara memperoleh kandidat terbaik jika ia mewawancarai semua
kandidat?  Berapakah peluangnya jika ia memilih kandidat pertama?

\item Anggaplah Barbara memutuskan untuk mewawancarai separuh pertama kandidat, lalu
melanjutkan wawancara sampai menemukan kandidat yang lebih baik daripada semua kandidat
yang telah ditemuinya.  Tunjukkan bahwa peluangnya untuk akhirnya memperoleh kandidat
terbaik lebih besar daripada 25 persen.
\end{enumerate}

\i\label{exer 3.1.24} Untuk tugas yang dijelaskan dalam Latihan~\ref{exer 3.1.23}, dapat
ditunjukkan\footnote{E.~B. Dynkin and A.~A. Yushkevich, \emx{Markov Processes: Theorems
and Problems,} trans.~J.~S. Wood (New York: Plenum, 1969).} bahwa strategi terbaik adalah
melewatkan $k - 1$ kandidat pertama, dengan $k$ sebagai
bilangan bulat terkecil yang memenuhi
$$
\frac 1k + \frac 1{k + 1} + \cdots + \frac 1{n - 1} \leq 1\ .
$$ Dengan strategi ini, peluang memperoleh kandidat terbaik kira-kira $1/e = .368$.
Tulislah program untuk menyimulasikan wawancara Barbara Smith jika ia menggunakan
strategi optimal ini dengan $n = 10$, lalu lihat apakah Anda dapat memverifikasi bahwa
peluang keberhasilannya kira-kira $1/e$.
\end{LJSItem}
\vspace{.25in}

\section{Kombinasi}\label{sec 3.2} 

Setelah menguasai permutasi, sekarang kita membahas kombinasi.  Misalkan $U$ suatu
himpunan dengan $n$ unsur; kita ingin menghitung banyaknya subhimpunan berbeda dari
himpunan $U$ yang memiliki tepat $j$ unsur.  Himpunan kosong dan himpunan $U$ dianggap
sebagai subhimpunan dari $U$.  Himpunan kosong biasanya dinyatakan dengan $\phi$.

\begin{example}\label{exam 3.7} Misalkan $U = \{a,b,c\}$.  Subhimpunan dari $U$ adalah
$$ \phi,\ \{a\},\ \{b\},\ \{c\},\ \{a,b\},\ \{a,c\},\ \{b,c\},\ \{a,b,c\}\ .$$ 
\end{example}

\subsection*{Koefisien Binomial}

Banyaknya subhimpunan berbeda dengan $j$ unsur yang dapat dipilih dari suatu himpunan
dengan $n$ unsur dinyatakan dengan ${n \choose j}$, dan dibaca ``$n$ pilih $j$."
Bilangan $n \choose j$ disebut \emx{koefisien binomial.}\index{koefisien binomial}
Istilah ini berasal dari penerapannya dalam aljabar yang akan dibahas kemudian pada
bagian ini.

Dalam contoh di atas, terdapat satu subhimpunan tanpa unsur, tiga subhimpunan dengan
tepat 1 unsur, tiga subhimpunan dengan tepat 2 unsur, dan satu subhimpunan dengan tepat
3 unsur.  Jadi, ${3 \choose 0} = 1$, ${3 \choose 1} = 3$, ${3 \choose 2} = 3$,
dan ${3 \choose 3} = 1$.   Perhatikan bahwa seluruhnya terdapat $2^3 = 8$ subhimpunan.
(Kita telah melihat bahwa suatu himpunan dengan $n$ unsur memiliki $2^n$ subhimpunan;
lihat Latihan~\ref{sec 3.1}.\ref{exer
3.1.8}.) Dengan demikian,

$$ {3 \choose 0} + {3 \choose 1}  + {3 \choose 2} + {3 \choose 3} = 2^3 = 8\ ,
$$
$$ {n \choose 0}  = {n \choose n} = 1\ .  
$$

Anggaplah $n > 0$.  Karena hanya ada satu cara untuk memilih himpunan tanpa unsur dan
hanya satu cara untuk memilih himpunan dengan $n$ unsur, nilai $n \choose j$ lainnya
ditentukan oleh \emx{relasi rekurensi} berikut:

\begin{theorem}\label{thm 3.6}                 
Untuk bilangan bulat $n$ dan $j$, dengan $0 < j < n$, koefisien binomial memenuhi:
\begin{equation} {n \choose j} = {{n-1} \choose j} + {{n-1} \choose {j - 1}}\ .
\label{eq 3.3}                                                                              
\end{equation}
\proof Kita ingin memilih subhimpunan dengan $j$ unsur.  Pilih sebuah unsur $u$ dari
$U$.  Pertama, anggaplah kita tidak menginginkan $u$ berada dalam subhimpunan.  Maka
kita harus memilih $j$ unsur dari himpunan dengan $n - 1$ unsur; hal ini dapat dilakukan
dengan ${{n-1} \choose j}$ cara.  Sebaliknya, anggaplah kita menginginkan $u$ berada
dalam subhimpunan.  Maka kita harus memilih $j - 1$ unsur lainnya dari $n - 1$ unsur
$U$ yang tersisa; hal ini dapat dilakukan dengan ${{n-1} \choose {j - 1}}$ cara.
Karena $u$ berada atau tidak berada dalam subhimpunan kita, banyaknya cara memilih
subhimpunan dengan $j$ unsur adalah jumlah banyaknya subhimpunan dengan $j$ unsur yang
memuat $u$ dan banyaknya yang tidak memuatnya---inilah yang dinyatakan oleh
Persamaan~\ref{eq 3.3}.
\end{theorem}

Koefisien binomial $n \choose j$ didefinisikan sebagai 0 jika $j < 0$ atau jika
$j > n$.  Dengan definisi ini, pembatasan pada $j$ dalam Teorema~\ref{thm 3.6} tidak
diperlukan.

\subsection*{Segitiga Pascal}

Relasi \ref{eq 3.3}, bersama dengan pengetahuan bahwa
$$ {n \choose 0} = {n \choose n }= 1\ ,
$$ menentukan sepenuhnya bilangan $n \choose j$.  Kita dapat menggunakan relasi ini
untuk menentukan \emx{segitiga Pascal}\index{segitiga Pascal} yang terkenal, yang
menampilkan semua bilangan tersebut dalam bentuk matriks (lihat Gambar~\ref{fig 3.6}).

\putfig{4.5truein}{PSfig3-6}{Segitiga Pascal.}{fig 3.6}   

Baris ke-$n$ segitiga ini memiliki entri $n \choose 0$,~$n \choose 1$,\dots,~$n
\choose n$.   Kita mengetahui bahwa bilangan pertama dan terakhir adalah 1.  Bilangan
lainnya ditentukan oleh relasi rekurensi Persamaan~\ref{eq 3.3}; yaitu, entri
${n \choose j}$ untuk $0 < j < n$ pada baris ke-$n$ segitiga Pascal merupakan
\emx{jumlah} entri yang tepat berada di atasnya dan entri yang tepat berada di sebelah
kirinya pada baris ke-$(n - 1)$.  Sebagai contoh, ${5 \choose 2} = 6 + 4 = 10$.

Algoritma untuk membangun segitiga Pascal ini dapat digunakan untuk menulis program
komputer yang menghitung koefisien binomial.  Anda diminta melakukannya dalam
Latihan~\ref{exer 3.2.4}.
\par
Meskipun segitiga Pascal menyediakan cara untuk membangun koefisien binomial secara
rekursif, kita juga dapat memberikan rumus untuk $n \choose j$.

\begin{theorem}\label{thm 3.7}  Koefisien binomial diberikan oleh rumus
\begin{equation} {n \choose j }= \frac{(n)_j}{j!}\ .
\label{eq 3.4}  
\end{equation}                                                                        
\proof Setiap subhimpunan berukuran $j$ dari suatu himpunan berukuran $n$ dapat
diurutkan dengan $j!$ cara.  Setiap pengurutan tersebut merupakan permutasi-$j$ dari
himpunan berukuran $n$.  Banyaknya permutasi-$j$ adalah $(n)_j$, sehingga banyaknya
subhimpunan berukuran $j$ adalah
$$
\frac{(n)_j}{j!}\ .
$$ Hal ini menyelesaikan bukti.
\end{theorem}

Rumus di atas dapat ditulis kembali dalam bentuk
$${n \choose j} = \frac{n!}{j!(n-j)!}\ .$$ Hal ini langsung menunjukkan bahwa
$$ {n \choose j} = {n \choose {n-j}}\ .
$$ Ketika menggunakan Persamaan~\ref{eq 3.4} untuk menghitung ${n \choose j}$, jika
perkalian dan pembagian dilakukan secara berselang-seling, semua nilai antara dalam
perhitungan merupakan bilangan bulat.  Selain itu, tidak satu pun nilai antara tersebut
melebihi nilai akhir.  (Lihat Latihan~\ref{exer 3.2.39}.)
\par
Hal lain yang perlu dinyatakan mengenai Persamaan~\ref{eq 3.4} adalah bahwa jika
persamaan tersebut digunakan untuk \emx {mendefinisikan} koefisien binomial, $n$ tidak
lagi harus berupa bilangan bulat positif.  Menurut definisi ini, peubah $j$ tetap harus
berupa bilangan bulat tak-negatif.  Gagasan ini berguna ketika memperluas Teorema
Binomial ke eksponen umum.  (Teorema Binomial untuk eksponen bilangan bulat tak-negatif
diberikan di bawah sebagai Teorema~\ref{thm 3.9}.)
\subsection*{Tangan Poker}\index{poker}

\begin{example}\label{exam 3.8} Pemain poker kadang-kadang bertanya-tanya mengapa
\emx{four of a kind} mengalahkan \emx{full house.}  Satu tangan poker merupakan
subhimpunan acak berisi 5 unsur dari setumpuk 52 kartu.  Suatu tangan memiliki four of
a kind jika memuat empat kartu bernilai sama---misalnya, empat kartu enam atau empat
raja.  Tangan itu merupakan full house jika memuat tiga kartu dari satu nilai dan dua
kartu dari nilai lain---misalnya, tiga kartu dua dan dua ratu.  Mari kita lihat tangan
mana yang lebih mungkin.  Berapa banyak tangan yang memiliki four of a kind?  Terdapat
13 cara untuk menentukan nilai keempat kartu.  Untuk setiap cara tersebut, terdapat
48 kemungkinan bagi kartu kelima.  Jadi, banyaknya tangan four of a kind adalah $13
\cdot 48 = 624$.  Karena jumlah seluruh tangan yang mungkin adalah ${52 \choose 5} =
2598960$, peluang tangan dengan four of a kind adalah $624/2598960 = .00024$.
\par
Sekarang perhatikan kasus full house; berapa banyak tangan seperti itu?  Terdapat 13
pilihan untuk nilai yang muncul tiga kali; untuk setiap pilihan tersebut, terdapat
${4 \choose 3} = 4$ pilihan bagi tiga kartu tertentu dengan nilai itu yang berada dalam
tangan.  Setelah memilih ketiga kartu tersebut, terdapat 12 kemungkinan untuk nilai
yang muncul dua kali; untuk setiap kemungkinan, terdapat ${4 \choose 2} = 6$ pilihan
bagi pasangan tertentu dengan nilai itu.  Jadi, banyaknya full house adalah $13 \cdot 4
\cdot 12 \cdot 6 = 3744$, dan peluang memperoleh tangan full house adalah
$3744/2598960 = .0014$.  Jadi, meskipun kedua jenis tangan jarang diperoleh, peluang
Anda memperoleh full house enam kali lebih besar daripada peluang memperoleh four of a kind.
\end{example}
\pagebreak[4]
\subsection*{Percobaan Bernoulli}

Penggunaan utama koefisien binomial akan muncul dalam kajian salah satu proses acak
penting yang disebut \emx{percobaan Bernoulli.}\index{proses percobaan Bernoulli}

\begin{definition}\label{def 3.4} Suatu \emx{proses percobaan Bernoulli} adalah urutan
$n$ percobaan acak sedemikian rupa sehingga
\begin{enumerate}
\item Setiap percobaan memiliki dua hasil yang mungkin, yang dapat kita sebut
\emx{keberhasilan} dan \emx{kegagalan.}
\item Peluang keberhasilan $p$ pada setiap percobaan sama untuk seluruh percobaan, dan
peluang ini tidak dipengaruhi oleh pengetahuan apa pun mengenai hasil sebelumnya.
Peluang kegagalan $q$ diberikan oleh $q = 1 - p$.
\end{enumerate}
\end{definition}

\putfig{4truein}{PSfig3-7}{Diagram pohon tiga percobaan Bernoulli.}{fig 3.7} 

\begin{example}\label{exam 3.9} Berikut adalah proses percobaan Bernoulli:
\begin{enumerate}
\item Sebuah koin dilempar sepuluh kali.  Dua hasil yang mungkin adalah sisi kepala dan
sisi ekor.  Peluang munculnya sisi kepala pada setiap lemparan adalah 1/2.
\item Suatu jajak pendapat dilakukan dengan menanyakan kepada 1000 orang yang dipilih
secara acak dari populasi apakah mereka mendukung Equal Rights Amendment---dua hasilnya
adalah ya dan tidak.  Peluang $p$ untuk jawaban ya (yaitu, keberhasilan) menunjukkan
proporsi orang dalam seluruh populasi yang mendukung amendemen tersebut.
\item Seorang penjudi memasang serangkaian taruhan 1 dolar, setiap kali bertaruh pada
warna hitam dalam rolet di Las~Vegas.  Di sini, keberhasilan berarti menang 1 dolar dan
kegagalan berarti kehilangan 1 dolar.  Karena dalam rolet Amerika penjudi menang jika
bola berhenti pada salah satu dari 18 di antara 38 posisi dan kalah jika tidak, peluang
menang adalah $p = 18/38 =
.474$.
\end{enumerate}
\end{example}

Untuk menganalisis proses percobaan Bernoulli, kita memilih pohon biner sebagai ruang
sampel dan menetapkan distribusi peluang pada lintasan-lintasan di pohon tersebut.
Sebagai contoh, misalkan kita memiliki tiga percobaan Bernoulli.  Hasil yang mungkin
ditunjukkan dalam diagram pohon pada Gambar~\ref{fig 3.7}.  Kita mendefinisikan $X$
sebagai peubah acak yang mewakili hasil proses, yaitu tripel terurut yang tersusun atas
simbol S dan F.  Peluang yang ditetapkan pada cabang pohon mewakili peluang setiap percobaan
tersendiri.  Misalkan hasil percobaan ke-$i$ dinyatakan dengan peubah acak $X_i$, dengan
fungsi distribusi $m_i$.  Karena kita telah menganggap bahwa hasil suatu percobaan tidak
memengaruhi hasil percobaan lainnya, kita menetapkan peluang yang sama pada setiap
tingkat pohon.  Suatu hasil $\omega$ untuk seluruh percobaan merupakan lintasan melalui
pohon.  Sebagai contoh, $\omega_3$ mewakili hasil SFS.  Penafsiran frekuensi atas
peluang membuat kita mengharapkan proporsi $p$ keberhasilan pada percobaan pertama; di
antaranya, proporsi $q$ kegagalan pada percobaan kedua; dan di antaranya lagi, proporsi
$p$ keberhasilan pada percobaan ketiga.  Hal ini menyarankan penetapan peluang $pqp$
pada hasil $\omega_3$.  Secara lebih umum, kita menetapkan fungsi distribusi
$m(\omega)$ untuk lintasan~$\omega$ dengan mendefinisikan $m(\omega)$ sebagai hasil kali
peluang cabang di sepanjang lintasan~$\omega$.  Jadi, peluang terjadinya tiga kejadian,
S pada percobaan pertama, F pada percobaan kedua, dan S pada percobaan ketiga, merupakan
hasil kali peluang masing-masing kejadian.  Pada bab berikutnya kita akan melihat bahwa
hal ini berarti kejadian-kejadian tersebut \emx{saling bebas}, dalam arti pengetahuan
tentang satu kejadian tidak memengaruhi prediksi kita mengenai terjadinya kejadian lain.

\subsection*{Peluang Binomial}

Kita terutama akan tertarik pada peluang bahwa dalam $n$ percobaan Bernoulli terdapat
tepat $j$ keberhasilan.  Peluang ini kita nyatakan dengan $b(n,p,j)$.  Mari kita hitung
nilai khusus $b(3,p,2)$ berdasarkan ukuran pohon kita.  Terlihat bahwa ada tiga lintasan yang
memiliki tepat dua keberhasilan dan satu kegagalan, yaitu $\omega_2$,~$\omega_3$,
dan~$\omega_5$.  Setiap lintasan ini memiliki peluang yang sama, $p^2q$.  Jadi,
$b(3,p,2) = 3p^2q$.  Dengan mempertimbangkan semua jumlah keberhasilan yang mungkin,
kita memperoleh

\begin{eqnarray*} 
b(3,p,0) &=& q^3\ ,\\
b(3,p,1) &=& 3pq^2\ ,\\ 
b(3,p,2) &=& 3p^2q\ ,\\
b(3,p,3) &=& p^3\ .
\end{eqnarray*}
Dengan cara yang sama, kita dapat membuat ukuran pohon untuk $n$ percobaan dan
menentukan $b(n,p,j)$ bagi kasus umum $n$ percobaan Bernoulli.

\begin{theorem}\label{thm 3.8} Diberikan $n$ percobaan Bernoulli dengan peluang
keberhasilan $p$ pada setiap percobaan, peluang tepat $j$ keberhasilan adalah
$$ b(n,p,j) = {n \choose j} p^j q^{n - j}
$$ dengan $q = 1 - p$.
\proof Kita membangun ukuran pohon sebagaimana dijelaskan di atas.  Kita ingin mencari
jumlah peluang semua lintasan yang memiliki tepat $j$ keberhasilan dan $n - j$
kegagalan.  Setiap lintasan tersebut diberi peluang $p^j q^{n - j}$.  Ada berapa banyak
lintasan seperti itu?  Untuk menentukan suatu lintasan, dari $n$ percobaan yang mungkin,
kita harus memilih suatu subhimpunan berisi $j$ percobaan sebagai keberhasilan, sedangkan
$n-j$ hasil yang tersisa merupakan kegagalan.  Hal ini dapat dilakukan dengan
$n \choose j$ cara.   Jadi, jumlah peluangnya adalah
$$ b(n,p,j) = {n \choose j} p^j q^{n - j}\ .
$$
\end{theorem}

\begin{example}\label{exam 3.10} Sebuah koin seimbang dilempar enam kali.  Berapakah
peluang muncul tepat tiga sisi kepala?  Jawabannya adalah
$$ b(6,.5,3) = {6 \choose 3} \left(\frac12\right)^3 \left(\frac12\right)^3 = 20 \cdot
\frac1{64} = .3125\ .
$$
\end{example}

\begin{example}\label{exam 3.11} Sebuah dadu dilempar empat kali.  Berapakah peluang
bahwa kita memperoleh tepat satu~6?  Kita memperlakukannya sebagai percobaan Bernoulli
dengan \emx{keberhasilan} = ``melempar 6" dan \emx{kegagalan} = ``melempar bilangan
selain 6."  Maka $p = 1/6$, dan peluang tepat satu keberhasilan dalam empat percobaan
adalah
$$ b(4,1/6,1) = {4 \choose 1 }\left(\frac16\right)^1 \left(\frac56\right)^3 = .386\ .
$$
\end{example}

Untuk menghitung peluang binomial dengan komputer, kalikan fungsi choose$(n,k)$ dengan
$p^kq^{n - k}$.  Program {\bf BinomialProbabilities}
\index{BinomialProbabilities (program)} mencetak peluang binomial $b(n, p, k)$ untuk $k$
di antara $kmin$ dan $kmax$, beserta jumlah peluang tersebut.  Kita telah menjalankan
program ini untuk $n = 100$, $p = 1/2$, $kmin = 45$, dan $kmax = 55$; keluarannya
ditunjukkan pada Tabel~\ref{table 3.27}.  Perhatikan bahwa peluang masing-masing cukup
kecil.  Peluang muncul tepat 50 sisi kepala dalam 100 lemparan koin adalah sekitar .08.
Intuisi kita mengatakan bahwa ini merupakan hasil yang paling mungkin, dan itu benar;
namun demikian, hasil tersebut bukanlah hasil yang sangat mungkin.

\begin{table}
\centering
\begin{tabular}{cc} $k$     & $b(n,p,k)$ \\
\\ 
45 & .0485 \\
46 & .0580 \\
47 & .0666 \\
48 & .0735 \\
49 & .0780 \\
50 & .0796 \\
51 & .0780 \\
52 & .0735 \\
53 & .0666 \\
54 & .0580 \\
55 & .0485 \\

\end{tabular}
\caption{Peluang binomial untuk $n = 100,\ p = 1/2$.}
\label{table 3.27}
\end{table}




\subsection*{Distribusi Binomial}

\begin{definition}\label{def 3.5} Misalkan $n$ bilangan bulat positif dan $p$ bilangan
real antara 0 dan 1.  Misalkan $B$ peubah acak yang menghitung jumlah keberhasilan dalam
proses percobaan Bernoulli dengan parameter $n$ dan $p$.  Maka distribusi $b(n, p, k)$
dari $B$ disebut \emx{distribusi binomial}.\index{distribusi binomial}
\index{fungsi distribusi!binomial}
\end{definition}

Kita dapat memperoleh gambaran yang lebih baik mengenai distribusi binomial dengan
membuat grafik distribusi ini untuk berbagai nilai $n$ dan $p$ (lihat
Gambar~\ref{fig 3.8}).  Plot dalam gambar tersebut dibuat menggunakan program
{\bf BinomialPlot}.\index{BinomialPlot (program)}

\putfig{4.5truein}{PSfig3-8}{Distribusi binomial.}{fig 3.8} 

\par Kita telah menjalankan program ini untuk $p = .5$ dan $p = .3$.  Perhatikan bahwa
bahkan untuk $p = .3$, grafiknya cukup simetris.  Kita akan memperoleh penjelasan
mengenai hal ini dalam Bab~\ref{chp 9}.   Kita juga memperhatikan bahwa peluang tertinggi
terjadi di sekitar nilai $np$, tetapi peluang tertinggi ini makin kecil ketika $n$
bertambah.  Dalam Bab~\ref{chp 6}, kita akan melihat bahwa $np$ merupakan nilai
\emx{rata-rata} atau nilai \emx{harapan} dari distribusi binomial $b(n,p,k)$.
\par
Contoh berikut memberikan cara yang baik untuk melihat distribusi binomial ketika $p =
1/2$.
\begin{example}\label{exam 3.2.1}
Sebuah \emx{papan Galton}\index{papan Galton} adalah papan tempat sejumlah besar butir
peluru BB dijatuhkan dari saluran di bagian atas papan dan dibelokkan oleh sejumlah pasak dalam
perjalanannya menuju bagian bawah papan.  Posisi akhir setiap peluru merupakan hasil
sejumlah pembelokan acak ke kiri atau ke kanan.  Kita telah menulis program
{\bf GaltonBoard}\index{GaltonBoard (program)} untuk menyimulasikan percobaan ini.
\par
Kita telah menjalankan program untuk kasus 20 baris pasak dan 10{,}000 peluru yang
dijatuhkan.  Hasil simulasi ini ditunjukkan pada Gambar~\ref{fig 2.22}.
\par
\putfig{4.5truein}{PSfig2-22}{Simulasi papan Galton.}{fig 2.22} 

Perhatikan bahwa jika kita menulis 0 setiap kali peluru dibelokkan ke kiri, dan 1 setiap
kali peluru dibelokkan ke kanan, lintasan peluru dapat dijelaskan dengan urutan 0 dan 1
sepanjang $n$, sama seperti pada pelemparan koin sebanyak $n$ kali.
\par
Distribusi yang ditunjukkan pada Gambar~\ref{fig 2.22} merupakan contoh distribusi
empiris, dalam arti distribusi itu muncul melalui serangkaian percobaan.  Seperti yang
diharapkan, distribusi empiris ini menyerupai distribusi binomial yang bersesuaian dengan
parameter $n = 20$ dan $p = 1/2$.
\end{example}



\subsection*{Pengujian Hipotesis}\index{pengujian hipotesis}

\begin{example}\label{exam 3.12} Misalkan aspirin biasa terbukti efektif mengatasi sakit
kepala pada 60~persen kasus, dan sebuah perusahaan obat mengklaim bahwa aspirin barunya
dengan zat tambahan khusus untuk sakit kepala lebih efektif.  Kita dapat menguji klaim
ini sebagai berikut: klaim mereka kita sebut \emx{hipotesis alternatif,} sedangkan
negasinya, yaitu bahwa zat tambahan itu tidak memiliki pengaruh yang berarti, kita sebut
\emx{hipotesis nol.}  Jadi, hipotesis nol menyatakan bahwa $p = .6$, sedangkan hipotesis
alternatif menyatakan bahwa $p > .6$, dengan $p$ sebagai peluang aspirin baru itu efektif.
\par Kita memberikan aspirin tersebut kepada $n$ orang untuk diminum ketika mereka sakit
kepala.  Kita ingin mencari suatu bilangan $m$, yang disebut \emx{nilai kritis} bagi
percobaan kita, sedemikian rupa sehingga kita menolak hipotesis nol jika setidaknya $m$
orang sembuh, dan menerimanya jika tidak.  Bagaimana kita menentukan nilai kritis ini?
\par Pertama, perhatikan bahwa kita dapat melakukan dua jenis galat.  Galat pertama,
yang dalam statistika sering disebut \emx{galat tipe~1}\index{galat tipe 1}, adalah
menolak hipotesis nol padahal sebenarnya benar.  Galat kedua, yang disebut \emx{galat
tipe~2,}\index{galat tipe 2} adalah menerima hipotesis nol ketika hipotesis itu salah.
Untuk menentukan peluang kedua jenis galat ini, kita memperkenalkan fungsi $\alpha(p)$,
yang didefinisikan sebagai peluang kita menolak hipotesis nol, dengan peluang tersebut
dihitung berdasarkan anggapan bahwa hipotesis nol benar.  Dalam kasus ini, kita memiliki
$$\alpha(p)  =  \sum_{m \leq k \leq n} b(n,p,k)\ .$$
\par 
Perhatikan bahwa $\alpha(.6)$ adalah peluang galat tipe~1, sebab ini merupakan peluang
jumlah keberhasilan yang tinggi bagi zat tambahan yang tidak efektif.  Jadi, untuk $n$
tertentu, kita ingin memilih $m$ agar $\alpha(.6)$ cukup kecil, sehingga kemungkinan
galat tipe~1 berkurang.  Namun, ketika $m$ meningkat melampaui nilai yang paling mungkin
$np
= .6n$, $\alpha(.6)$, sebagai ekor atas distribusi binomial, mendekati 0.  Jadi,
\emx{meningkatkan} $m$ membuat galat tipe~1 kurang mungkin terjadi.
\par Sekarang misalkan zat tambahan itu memang efektif, sehingga $p$ jauh lebih besar
daripada .6; katakanlah $p = .8$.  (Nilai alternatif $p$ ini dipilih secara sembarang;
perhitungan berikut bergantung pada pilihan tersebut.)  Memilih $m$ jauh di bawah $np =
.8n$ akan meningkatkan $\alpha(.8)$, sebab sekarang $\alpha(.8)$ mencakup hampir seluruh
distribusi binomial kecuali ekor bawahnya.  Bahkan, jika kita menetapkan
$\beta(.8) = 1 - \alpha(.8)$, maka $\beta(.8)$ memberi kita peluang galat tipe~2,
sehingga \emx{menurunkan} $m$ membuat galat tipe~2 kurang mungkin terjadi.
\par Produsen ingin menghindari galat tipe~2, sebab jika galat tersebut terjadi,
pengujian tidak menunjukkan bahwa obat baru lebih baik, padahal sebenarnya demikian.
Jika nilai alternatif $p$ dipilih lebih dekat dengan nilai $p$ dalam hipotesis nol
(dalam kasus ini $p =.6$), maka untuk populasi uji tertentu, nilai $\beta$ akan meningkat.
Jadi, jika ahli statistik produsen memilih nilai alternatif $p$ yang dekat dengan nilai
dalam hipotesis nol, diperlukan biaya besar (yaitu, populasi uji harus besar) untuk
menolak hipotesis nol dengan nilai $\beta$ yang kecil.
\par Jadi, untuk populasi uji $n$ tertentu, kita berharap dapat memilih nilai $m$ yang,
jika mungkin, membuat kedua peluang ini kecil.  Jika melakukan galat tipe~1, kita akan
membeli banyak aspirin yang pada dasarnya biasa dengan harga yang terlalu tinggi;
galat tipe~2 berarti kita melewatkan kesempatan memperoleh obat yang lebih unggul dengan
harga murah.  Katakanlah kita ingin bilangan kritis $m$ membuat peluang masing-masing
kasus yang tidak diinginkan ini kurang dari 5 persen.
\par Kita menulis program {\bf PowerCurve}\index{PowerCurve (program)} untuk memplot,
bagi $n = 100$ dan beberapa nilai $m$ yang dipilih, fungsi $\alpha(p)$, dengan $p$
berkisar dari .4 hingga~1.  Hasilnya ditunjukkan pada Gambar~\ref{fig 3.9}.  Dalam grafik,
kita menyertakan sebuah kotak (dengan garis putus-putus) dari .6 hingga .8, dengan sisi
bawah dan atas pada ketinggian .05 dan .95.  Suatu nilai $m$ memenuhi persyaratan kita
jika dan hanya jika grafik $\alpha$ memasuki kotak dari bawah dan keluar dari atas
(mengapa?---mana kriteria tipe~1 dan mana kriteria tipe~2?).  Ketika $m$ meningkat,
grafik $\alpha$ bergeser ke kanan.  Beberapa percobaan menunjukkan bahwa $m = 69$
merupakan nilai terkecil $m$ yang mencegah galat tipe~1, sedangkan $m = 73$ merupakan
nilai terbesar yang mencegah galat tipe~2.  Jadi, kita dapat memilih nilai kritis antara
69 dan 73.  Jika kita lebih ingin menghindari galat tipe~1, kita memilih 73; demikian
pula, kita memilih 69 jika menganggap galat tipe~2 lebih buruk.  Tentu saja, perusahaan
obat mungkin tidak senang dengan peluang galat sebesar 5 persen.  Mereka mungkin
menuntut peluang galat sebesar 1 persen.  Untuk itu, kita harus meningkatkan jumlah
percobaan $n$ (lihat Latihan~\ref{exer 3.2.28}).
\end{example}

\putfig{4.5truein}{PSfig3-9}{Kurva kuasa.}{fig 3.9} 

\subsection*{Ekspansi Binomial}

Selanjutnya kita mengingatkan pembaca tentang penerapan koefisien binomial dalam aljabar.
Penerapan ini adalah \emx{ekspansi binomial,} yang menjadi asal istilah koefisien binomial.
\begin{theorem}\label{thm 3.9}{\bf (Teorema Binomial)}\index{Teorema Binomial} Kuantitas $(a +
b)^n$ dapat dinyatakan dalam bentuk
$$ (a + b)^n = \sum_{j = 0}^n {n \choose j} a^j b^{n - j}\ .
$$
\proof Untuk melihat bahwa ekspansi ini benar, tuliskan
$$ (a + b)^n = (a + b)(a + b) \cdots (a + b)\ .
$$ Ketika mengalikannya, kita akan memperoleh jumlah suku yang masing-masing dihasilkan
dari pilihan sebuah $a$ atau $b$ untuk setiap faktor dari $n$ faktor tersebut.  Ketika kita memilih
$j$ buah $a$ dan $(n - j)$ buah $b$, kita memperoleh suku berbentuk $a^j b^{n - j}$.
Untuk menentukan suku tersebut, kita harus menetapkan $j$ dari $n$ suku dalam hasil
kali yang menjadi tempat kita memilih $a$.  Hal ini dapat dilakukan dengan
$n \choose j$ cara.   Jadi, mengumpulkan suku-suku ini dalam penjumlahan menghasilkan
suku ${n \choose j} a^j b^{n - j}$.
\end{theorem}
\par
Sebagai contoh, kita memiliki
\begin{eqnarray*} (a + b)^0 & = & 1 \\ (a + b)^1 & = & a + b \\ (a + b)^2 & = & a^2 +
2ab + b^2 \\ (a + b)^3 & = & a^3 + 3a^2b + 3ab^2 + b^3\ .
\end{eqnarray*}
Di sini terlihat bahwa koefisien pangkat-pangkat berurutan memang menghasilkan segitiga
Pascal.\index{segitiga Pascal}

\begin{corollary}\label{cor 3.1}  Jumlah unsur pada baris ke-$n$ segitiga Pascal adalah
$2^n$.  Jika unsur pada baris ke-$n$ segitiga Pascal dijumlahkan dengan tanda
berselang-seling, jumlahnya adalah~0.
\proof Pernyataan pertama dalam akibat ini mengikuti fakta bahwa
$$ 2^n = (1 + 1)^n = {n \choose 0} + {n \choose 1} + {n \choose 2} + \cdots + {n
\choose n}\ ,
$$ dan pernyataan kedua mengikuti fakta bahwa
$$ 0 = (1 - 1)^n = {n \choose 0} - {n \choose 1} + {n \choose 2}- \cdots + {(-1)^n}{n
\choose n}\ .
$$
\end{corollary}

Pernyataan pertama dari akibat tersebut memberi tahu kita bahwa banyaknya subhimpunan
dari suatu himpunan dengan $n$ unsur adalah $2^n$.  Kita akan menggunakan pernyataan
kedua dalam penerapan teorema binomial berikutnya.

Kita telah melihat bahwa ketika $A$ dan $B$ adalah sembarang dua kejadian
(bdk.~Bagian~\ref{sec 1.2}),
$$ P(A \cup B) = P(A) + P(B) - P(A \cap B).
$$ Sekarang kita memperluas teorema ini menjadi versi yang lebih umum, yang memungkinkan
kita mencari peluang bahwa setidaknya satu dari sejumlah kejadian terjadi.

\subsection*{Prinsip Inklusi-Eksklusi}\index{Prinsip Inklusi-Eksklusi}

\begin{theorem}\label{thm 3.10} Misalkan $P$ suatu distribusi peluang pada ruang sampel
$\Omega$, dan misalkan $\{A_1,\ A_2,\ \dots,\ A_n\}$ suatu himpunan kejadian berhingga.
Maka
$$ P(A_1 \cup A_2 \cup \cdots \cup A_n)  =  \sum_{i = 1}^n P(A_i)\ - 
\sum_{1 \leq i < j \leq n} P(A_i \cap A_j) 
$$
\begin{equation} 
\ \ \ \ \ \ \ \ \ \ \ \ \ \ \ \ \ \ \ \ \ \ \ \ \ \ \ \ \ \ \ \ \ \ \ \ \ \ \ \ \ \ \
\ 
\ \ \ \ \ \ \ \ \ \ \ \ \  + \sum_{1 \leq i < j < k \leq n} P(A_i \cap A_j \cap A_k) -
\cdots\ .\label{eq 3.5}
\end{equation}
Artinya, untuk mencari peluang bahwa setidaknya satu dari $n$ kejadian $A_i$ terjadi,
mula-mula jumlahkan peluang setiap kejadian, kemudian kurangi dengan peluang semua irisan
pasangan kejadian, tambahkan peluang semua irisan tiga kejadian, dan seterusnya.
\proof Jika hasil $\omega$ terjadi dalam setidaknya satu kejadian $A_i$, peluangnya
ditambahkan tepat sekali oleh ruas kiri Persamaan~\ref{eq 3.5}.  Kita harus menunjukkan
bahwa peluang itu ditambahkan tepat sekali oleh ruas kanan Persamaan~\ref{eq 3.5}.
Anggaplah $\omega$ berada dalam tepat $k$ himpunan.   Maka peluangnya ditambahkan $k$
kali dalam suku pertama, dikurangkan $k \choose 2$ kali dalam suku kedua, ditambahkan
$k \choose 3$ kali dalam suku ketiga, dan seterusnya.   Jadi, jumlah keseluruhan peluang
itu ditambahkan adalah
$$ {k \choose 1} - {k \choose 2} + {k \choose 3} - \cdots {(-1)^{k-1}} {k \choose k}\ .
$$ Namun,
$$ 0 = (1 - 1)^k = \sum_{j = 0}^k {k \choose j} (-1)^j = {k \choose 0} - \sum_{j = 1}^k
{k \choose j} {(-1)^{j - 1}}\ .
$$ Oleh karena itu,
$$ 1 = {k \choose 0} = \sum_{j = 1}^k {k \choose j} {(-1)^{j - 1}}\ .
$$
Jika hasil $\omega$ tidak berada dalam satu pun kejadian $A_i$, hasil tersebut tidak
dihitung pada kedua ruas persamaan.
\end{theorem}

\subsection*{Masalah Penitipan Topi}

\begin{example}\label{exam 3.13}  Kita kembali ke masalah penitipan topi
\index{masalah penitipan topi} yang dibahas dalam Bagian~\ref{sec 3.1}, yaitu masalah
mencari peluang bahwa suatu permutasi acak memuat setidaknya satu titik tetap.  Ingatlah
bahwa permutasi merupakan pemetaan satu-ke-satu dari suatu himpunan
$A = \{a_1,a_2,\dots,a_n\}$ ke dirinya sendiri.  Misalkan $A_i$ adalah kejadian bahwa
unsur ke-$i$, $a_i$, tetap pada pemetaan tersebut.  Jika kita mengharuskan $a_i$ tetap,
pemetaan atas $n - 1$ unsur lainnya memberikan sembarang permutasi dari $(n -
1)$ objek.  Karena terdapat $(n - 1)!$ permutasi demikian,
$P(A_i) = (n - 1)!/n! =
1/n$.  Karena terdapat $n$ pilihan untuk $a_i$, suku pertama
Persamaan~\ref{eq 3.5} adalah 1.  Dengan cara yang sama, agar suatu pasangan tertentu
$(a_i,a_j)$ tetap, kita dapat memilih sembarang permutasi atas $n - 2$ unsur yang
tersisa; terdapat $(n - 2)!$ pilihan demikian, sehingga
$$ P(A_i \cap A_j) = \frac{(n - 2)!}{n!} = \frac 1{n(n - 1)}\ .
$$ Banyaknya suku berbentuk demikian di ruas kanan Persamaan~\ref{eq 3.5} adalah $$ {n
\choose 2} = \frac{n(n - 1)}{2!}\ .
$$ Jadi, suku kedua Persamaan~\ref{eq 3.5} adalah
$$ -\frac{n(n - 1)}{2!} \cdot \frac 1{n(n - 1)} = -\frac 1{2!}\ .
$$ Dengan cara yang sama, untuk sembarang tiga kejadian tertentu $A_i$, $A_j$, $A_k$,
$$ P(A_i \cap A_j \cap A_k) = \frac{(n - 3)!}{n!} = \frac 1{n(n - 1)(n - 2)}\ ,
$$ dan banyaknya suku demikian adalah
$$ {n \choose 3} = \frac{n(n - 1)(n - 2)}{3!}\ ,
$$ sehingga suku ketiga Persamaan~\ref{eq 3.5} sama dengan 1/3!.  Dengan melanjutkan
cara ini, kita memperoleh
$$ P(\mbox {paling\ sedikit\ satu\ titik\ tetap}) = 1 - \frac 1{2!} + \frac 1{3!} - \cdots
(-1)^{n-1} \frac 1{n!}
$$ dan
$$ P(\mbox {tanpa\ titik\ tetap}) = \frac 1{2!} - \frac 1{3!} + \cdots (-1)^n \frac 1{n!}\ .
$$

Dari kalkulus kita mengetahui bahwa
$$ e^x = 1 + x + \frac 1{2!}x^2 + \frac 1{3!}x^3 + \cdots + \frac 1{n!}x^n + \cdots\ .
$$ Jadi, jika $x = -1$, kita memiliki
\begin{eqnarray*} e^{-1} & = &\frac 1{2!} - \frac 1{3!} + \cdots + \frac{(-1)^n}{n!} +
\cdots \\
       & = & .3678794\ .
\end{eqnarray*} Oleh karena itu, peluang tidak adanya titik tetap, yaitu tidak seorang
pun dari $n$ orang memperoleh kembali topinya sendiri, sama dengan jumlah $n$ suku
pertama dalam ungkapan untuk $e^{-1}$.  Deret ini konvergen sangat cepat.  Menghitung
jumlah parsial untuk $n = 3$ hingga~10 memberikan data pada Tabel~\ref{table 3.7}.
\begin{table}
\centering
\begin{tabular}{rl}
      &        Peluang tidak seorang pun  \\
   n  &        memperoleh topinya sendiri \\\hline
   3  &       \hspace{.25in}.333333     \\
   4  &       \hspace{.2in} .375        \\
   5  &       \hspace{.25in}.366667     \\
   6  &       \hspace{.25in}.368056     \\
   7  &       \hspace{.25in}.367857     \\
   8  &       \hspace{.25in}.367882     \\
   9  &       \hspace{.25in}.367879     \\
  10  &       \hspace{.25in}.367879     \\\hline
\end{tabular}
\caption{Masalah penitipan topi.}
\label{table 3.7}
\end{table}
\par
Setelah $n = 9$, peluang-peluang tersebut pada dasarnya sama hingga enam angka penting.
Menariknya, peluang tidak adanya titik tetap meningkat dan menurun secara bergantian
ketika $n$ bertambah.  Terakhir, kita mencatat bahwa hasil eksak kita sangat sesuai
dengan simulasi yang dilaporkan pada bagian sebelumnya.
\end{example}

\subsection*{Memilih Ruang Sampel}

Sekarang kita memiliki beberapa alat yang diperlukan untuk mendeskripsikan ruang sampel
secara akurat dan menetapkan fungsi peluang pada ruang sampel tersebut.  Meskipun
demikian, dalam beberapa kasus, proses deskripsi dan penetapan ini agak arbitrer.
Tentu saja, kita berharap ruang sampel yang dipilih beserta fungsi peluang yang ditetapkan
padanya menghasilkan model yang secara akurat memprediksi apa yang akan terjadi jika
percobaan benar-benar dilakukan.  Seperti ditunjukkan contoh-contoh berikut, ada keadaan
ketika deskripsi ruang sampel yang ``masuk akal" tidak menghasilkan model yang sesuai dengan
data.

Dalam buku Feller,\index{FELLER, W.}\footnote{W.~Feller, \emx{Introduction to Probability Theory and
Its Applications} vol.~1, 3rd~ed. (New York: John Wiley and Sons, 1968), p. 41}
diberikan sepasang model yang menjelaskan susunan jenis-jenis partikel elementer tertentu,
seperti foton\index{foton} dan proton\index{proton}.  Eksperimen ternyata menunjukkan
bahwa jenis partikel elementer tertentu memperlihatkan perilaku yang dijelaskan secara
akurat oleh satu model, yang disebut \emx{``statistika Bose-Einstein,"}
\index{statistika Bose-Einstein}, sedangkan jenis partikel elementer lainnya dapat
dimodelkan menggunakan \emx{``statistika Fermi-Dirac."}\index{statistika Fermi-Dirac}
Feller\index{FELLER, W.} menggunakan contoh ini untuk menegaskan bahwa pemilihan model
peluang memerlukan bukti empiris: argumen \emx{a priori} saja tidak dapat menentukan bahwa
foton dan proton harus mengikuti hukum peluang yang berbeda. Ini merupakan ringkasan editorial;
ungkapan sumber tidak diterjemahkan atau direproduksi.

Sekarang kita memberikan beberapa contoh proses penjelasan dan penetapan ini.

\begin{example}\label{exam 3.14}
 Dalam model mekanika kuantum\index{mekanika kuantum} untuk atom helium\index{helium},
berbagai parameter dapat digunakan untuk mengklasifikasikan keadaan energi atom.  Dalam
keadaan spin triplet ($S =
1$) dengan momentum sudut orbital 1 ($L = 1$), terdapat tiga kemungkinan, 0, 1, atau 2,
untuk momentum sudut total ($J$).  (Pembaca tidak diharapkan memahami arti semua ini;
bahkan, contoh tersebut lebih jelas jika pembaca \emx{tidak} mengetahui apa pun tentang
mekanika kuantum.)  Kita ingin menetapkan peluang pada tiga kemungkinan bagi $J$.
Anda tentu akan menolak gagasan untuk menetapkan peluang $1/3$ pada setiap hasil
tersebut.  Sekarang tanyakan kepada diri sendiri mengapa Anda menolak penetapan ini.
Jawabannya mungkin karena Anda tidak memiliki ``intuisi" (yaitu,
pengalaman) mengenai perilaku atom helium.  Sebenarnya, dalam contoh ini, peluang
$1/9,\ 3/9,$ dan $5/9$ ditetapkan oleh teori.  Teori memberikan penetapan ini karena
frekuensi tersebut diamati \emx{dalam eksperimen}, lalu parameter lebih lanjut
dikembangkan dalam teori agar frekuensi tersebut dapat diprediksi.
\end{example}

\begin{example}\label{exam 3.15}  Misalkan dua keping uang logam satu sen masing-masing
dilempar satu kali.   Ada beberapa cara yang ``masuk akal" untuk mendeskripsikan ruang sampel.  Salah satu
caranya adalah menghitung jumlah sisi kepala dalam hasil; dalam kasus ini, ruang sampel
dapat ditulis $\{0, 1,
2\}$.  Deskripsi lain atas ruang sampel adalah himpunan semua pasangan terurut $H$ dan
$T$,
yaitu \[ \{(H,H), (H, T), (T, H), (T, T)\}. \]
\noindent Kedua deskripsi ini akurat, tetapi mudah dilihat bahwa (paling banyak) hanya
salah satunya, jika diberi fungsi peluang konstan, yang dapat mengklaim memodelkan
kenyataan secara akurat.  Dalam kasus ini, tidak seperti contoh sebelumnya, pembaca
mungkin akan mengatakan bahwa deskripsi kedua, dengan setiap hasil diberi peluang
$1/4$, merupakan deskripsi yang ``benar".  Keyakinan ini berasal dari pengalaman;
tidak ada bukti bahwa kenyataan bekerja dengan cara demikian.
\end{example} 
\par
Pembaca juga diarahkan ke Latihan~\ref{exer 3.2.26} untuk contoh lain dari proses ini.

\subsection*{Catatan Historis}

Koefisien binomial memiliki sejarah panjang dan penuh warna yang berpuncak pada karya
Pascal\index{PASCAL, B.}, \emx{Treatise on the Arithmetical Triangle,}
\footnote{B.~Pascal, \emx{Trait\'e du Triangle Arithm\'etique} (Paris: Desprez, 1665).}
tempat Pascal mengembangkan banyak sifat penting bilangan-bilangan ini.  Sejarah tersebut
dipaparkan dalam buku \emx{Pascal's Arithmetical Triangle} karya A.~W.~F. Edwards.
\index{EDWARDS, A. W. F.}\footnote{A.~W.~F. Edwards, \emx{Pascal's Arithmetical Triangle}
(London: Griffin, 1987).}  Pascal menulis segitiganya\index{segitiga Pascal} dalam bentuk
yang ditunjukkan pada Tabel~\ref{table 3.8}.
\begin{table}
\centering
\begin{tabular}{llrrrrrrll} 
1  & 1 & 1  & 1  & 1  & 1   & 1  & 1  & 1 & 1\cr 
1  & 2 & 3  & 4  & 5  & 6   & 7  & 8  & 9  \cr  
1  & 3 & 6  & 10 & 15 & 21  & 28 & 36 \cr  
1  & 4 & 10 & 20 & 35 & 56  & 84 \cr  
1  & 5 & 15 & 35 & 70 & 126 \cr  
1  & 6 & 21 & 56 &126 \cr  
1  & 7 & 28 & 84 \cr  
1  & 8 & 36 \cr  
1  & 9 \cr  
1 \cr
\end{tabular}
\caption{Segitiga Pascal.}
\label{table 3.8}
\end{table}

\par Edwards menelusuri tiga cara berbeda munculnya koefisien binomial.  Ia menyebutnya
\emx{bilangan figuratif,} \emx{bilangan kombinatorial,} dan \emx{bilangan binomial.}
Semuanya merupakan nama untuk hal yang sama (yang kita sebut koefisien binomial), tetapi
kesamaan ketiganya baru disadari pada abad keenam belas.

\emx{Bilangan figuratif}\index{bilangan figuratif} berawal dari ketertarikan kaum
Pythagoras pada pola bilangan sekitar 540~{\footnotesize{SM}.}   Kaum Pythagoras
mempertimbangkan, misalnya, pola segitiga yang ditunjukkan pada Gambar~\ref{fig 3.10}.
Barisan bilangan
$$ 1, 3, 6, 10, \dots
$$ yang diperoleh sebagai banyaknya titik dalam setiap segitiga disebut \emx{bilangan
segitiga.}\index{bilangan segitiga}  Dari segitiga tersebut jelas bahwa bilangan
segitiga ke-$n$ hanyalah jumlah $n$ bilangan bulat pertama.  \emx{Bilangan tetrahedral}
\index{bilangan tetrahedral} merupakan jumlah bilangan segitiga dan diperoleh oleh
matematikawan Yunani Theon dan Nicomachus pada awal abad kedua~{\footnotesize{SM}.}
Sebagai contoh, bilangan tetrahedral 10 memiliki representasi geometris yang ditunjukkan
pada Gambar~\ref{fig 3.11}.  Tiga jenis pertama bilangan figuratif dapat ditampilkan
dalam bentuk tabel seperti pada Tabel~\ref{table 3.9}.

\putfig{4.5truein}{PSfig3-10}{Pola segitiga Pythagoras.}{fig 3.10} 

\putfig{2truein}{PSfig3-11}{Representasi geometris bilangan tetrahedral 10.}{fig
3.11} 


\begin{table}
\centering
$$
\begin{array}{llllllllll} 
 \mbox {bilangan\ asli}       & 1\hskip.2in & 2\hskip.2in & 3\hskip.2in & 4\hskip.2in & 5\hskip.2in & 6\hskip.2in & 7\hskip.2in & 8\hskip.2in & 9  \cr
 \mbox {bilangan\ segitiga}   & 1     & 3     & 6    & 10   & 15   & 21   & 28   & 36  & 45 \cr
 \mbox {bilangan\ tetrahedral}& 1     & 4     & 10   & 20   & 35   & 56   & 84   & 120 &165  
\end{array}
$$
\caption{Bilangan figuratif.}
\label{table 3.9}
\end{table}
\par Bilangan-bilangan ini memberikan empat baris pertama segitiga Pascal, tetapi tabel
tersebut baru diselesaikan di Barat pada abad keenam belas.
\par
Di Timur, matematikawan Hindu mulai menemukan koefisien binomial dalam masalah
kombinatorial.  Dalam \emx{Lilavati} dari 1150, Bhaskara memberikan kaidah untuk mencari
banyaknya sediaan obat yang menggunakan 1, 2, 3, 4, 5, atau~6 bahan yang mungkin.
\footnote{ibid., p.~27.}  Kaidahnya setara dengan rumus kita
$$ {n \choose r} = \frac{(n)_r}{r!}\ .
$$
\begin{figure}[ht]
\centering
\small
\begin{tabular}{c}
$1$\\
$1\quad 1$\\
$1\quad 2\quad 1$\\
$1\quad 3\quad 3\quad 1$\\
$1\quad 4\quad 6\quad 4\quad 1$\\
$1\quad 5\quad 10\quad 10\quad 5\quad 1$\\
$1\quad 6\quad 15\quad 20\quad 15\quad 6\quad 1$\\
$1\quad 7\quad 21\quad 34\quad 35\quad 21\quad 7\quad 1$\\
$1\quad 8\quad 28\quad 56\quad 70\quad 56\quad 28\quad 8\quad 1$
\end{tabular}
\caption{Transkripsi numerik mandiri susunan segitiga Chu Shih-chieh hingga
pangkat kedelapan; kesalahan penyalinan historis yang dibahas di bawah
dipertahankan.}
\label{fig 3.12}
\end{figure}
\par
Bilangan binomial sebagai koefisien dari $(a + b)^n$ muncul dalam karya matematikawan
di Tiongkok sekitar 1100.  Sumber-sumber dari sekitar masa ini menyebut adanya metode
tabulasi untuk memperoleh koefisien binomial.  Segitiga yang memberikan koefisien hingga
pangkat kedelapan disajikan oleh Chu Shih-chieh\index{CHU, S.-C.} dalam sebuah buku yang
ditulis sekitar 1303. Gambar~\ref{fig 3.12} menyajikan kembali susunan
matematisnya, bukan faksimile sumber historis.\footnote{Rujukan historis: J.\
Needham, \emx{Science and Civilization in China,} vol.~3 (New York: Cambridge
University Press, 1959), p.~135.}
Naskah asli buku Chu telah hilang, tetapi salinannya masih bertahan.  Edwards mencatat
bahwa ada kesalahan dalam salinan segitiga Chu ini.  Dapatkah Anda menemukannya?
(\emx {Petunjuk}: Dua bilangan yang seharusnya sama ternyata tidak sama.)  Salinan lain
tidak memuat kesalahan ini.
\par
Kemunculan pertama segitiga Pascal di Barat tampaknya berasal dari perhitungan Tartaglia
tentang banyaknya kemungkinan hasil lemparan $n$ dadu.\index{TARTAGLIA, N.}
\footnote{N. Tartaglia, \emx{General Trattato di Numeri et Misure} (Vinegia, 1556).}
Untuk satu dadu, jawabannya jelas 6.  Untuk dua dadu, kemungkinan hasilnya dapat
ditampilkan seperti pada Tabel~\ref{table 3.10}.

\begin{table}
\centering
$$\matrix{  11\cr   12 & 22\cr   13 & 23 & 33\cr   14 & 24 & 34 & 44\cr   15 & 25 & 35
& 45 & 55\cr   16 & 26 & 36 & 46 & 56 & 66\ \cr }
$$
\caption{Hasil pelemparan dua dadu.}
\label{table 3.10}
\end{table}

Menampilkannya dengan cara ini menyarankan bilangan segitiga keenam
$1 + 2 + 3 + 4 + 5 + 6 =
21$ untuk pelemparan 2 dadu.  Menurut catatan Tartaglia, ia memikirkan masalah itu
semalaman di Verona pada hari pertama Prapaskah 1523.
\footnote{Ringkasan berdasarkan Edwards, op.\ cit., p.~37; ungkapan sumber tidak
diterjemahkan atau direproduksi.} Ia menyadari bahwa perluasan tabel
figuratif memberikan jawaban untuk $n$ dadu.  Masalah tersebut terpikir oleh Tartaglia
setelah melihat orang meramal horoskop mereka sendiri menggunakan sebuah \emx{ Book of
Fortune,} untuk memilih syair melalui proses yang mencakup pencatatan angka pada muka
tiga dadu.  Sebanyak 56 kemungkinan hasil tiga dadu dicantumkan pada setiap halaman.
Cara bilangan ditulis dalam buku itu tidak menunjukkan kaitannya dengan bilangan
figuratif, tetapi metode pencacahan yang serupa dengan metode kita untuk 2 dadu
menunjukkannya.  Tabel Tartaglia baru diterbitkan pada 1556.
\par
Sebuah tabel koefisien binomial diterbitkan pada 1554 oleh matematikawan Jerman, Stifel.
\index{STIFEL, M.}\footnote{M. Stifel, \emx{ Arithmetica Integra} (Norimburgae, 1544).}
Segitiga Pascal juga muncul dalam \emx{Opus novum} karya Cardano dari 1570.
\index{CARDANO, G.}\footnote{G. Cardano, \emx{Opus Novum de Proportionibus Numerorum}
(Basilea, 1570).}  Cardano tertarik pada masalah mencari banyaknya cara memilih $r$
objek dari $n$.  Jadi, pada masa karya Pascal, segitiganya telah muncul sebagai hasil
kajian bilangan figuratif, bilangan kombinatorial, dan bilangan binomial, dan fakta bahwa
ketiganya sama tampaknya telah dipahami dengan cukup baik.
\par
Ketertarikan Pascal\index{PASCAL, B.|(}\index{FERMAT, P.|(} pada bilangan binomial
berasal dari surat-menyuratnya dengan Fermat mengenai masalah yang dikenal sebagai
masalah pembagian taruhan.\index{masalah pembagian taruhan}  Masalah ini dan
surat-menyurat antara Pascal dan Fermat dibahas dalam Bab~\ref{chp 1}.  Pembaca akan
mengingat bahwa masalah ini dapat dijelaskan sebagai berikut: Dua pemain A dan B
memainkan serangkaian permainan, dan pemain pertama yang memenangkan $n$ permainan
memenangkan pertandingan.  Kita ingin mencari peluang A memenangkan pertandingan pada
saat A telah memenangkan $a$ permainan dan B telah memenangkan $b$ permainan.  (Lihat
Latihan~\ref{sec 4.1}.\ref{exer 5.1.11}-\ref{sec
4.1}.\ref{exer 5.1.13}.)
\par
Pascal menyelesaikan masalah tersebut dengan induksi mundur, hampir sama dengan cara
kita menulis program komputer untuk menyelesaikannya sekarang.  Ia merujuk pada metode
kombinatorial Fermat yang berjalan sebagai berikut: Jika A memerlukan $c$ kemenangan dan
B memerlukan $d$ kemenangan untuk menang, kita meminta para pemain terus bermain sampai
mereka memainkan $c + d - 1$ permainan.  Pemenang rangkaian yang diperpanjang ini akan
sama dengan pemenang rangkaian asli.  Peluang A menang dalam rangkaian yang diperpanjang,
dan dengan demikian dalam rangkaian asli, adalah
$$
\sum_{r = c}^{c + d - 1} \frac 1{2^{c + d - 1}} {{c + d - 1} \choose r}\ .
$$ 
Bahkan pada saat surat-surat itu ditulis, Pascal tampaknya telah memahami rumus ini.
\par Misalkan pemain pertama yang memenangkan $n$ permainan memenangkan pertandingan,
dan misalkan setiap pemain memasang taruhan sebesar $x$.  Pascal mempelajari nilai
kemenangan dalam suatu permainan tertentu.  Yang dimaksudnya adalah peningkatan nilai
harapan perolehan bagi pemenang permainan tertentu yang sedang dibahas.  Ia menunjukkan
bahwa nilai permainan pertama adalah
$$ \frac {1\cdot3\cdot5\cdot\dots\cdot(2n - 1)}{2\cdot4\cdot6\cdot\dots\cdot(2n)}x\ .
$$ 
Buktinya tampaknya menggunakan rumus Fermat dan fakta bahwa rasio hasil kali bilangan
ganjil terhadap hasil kali bilangan genap di atas sama dengan peluang muncul tepat $n$
sisi kepala dalam $2n$ lemparan koin.  (Lihat Latihan~\ref{exer 3.2.38}.)
\par
Pascal menyajikan kepada Fermat tabel yang ditunjukkan pada Tabel~\ref{table 3.11}.
\begin{table}
\centering
\resizebox{\textwidth}{!}{%
\begin{tabular}{ccccccc}\hline
                          &\multicolumn{6}{c}{jika masing-masing mempertaruhkan 256 dalam} \\
Dari 256 milik lawan saya &   6    & 5     & 4     & 3     & 2     & 1     \\ 
yang saya peroleh, untuk  & permainan & permainan & permainan & permainan & permainan & permainan \\ \cline{2-7}
permainan ke-1            & 63     & 70    & 80    & 96    & 128   & 256   \\
permainan ke-2            & 63     & 70    & 80    & 96    & 128   \\
permainan ke-3            & 56     & 60    & 64    & 64\\
permainan ke-4            & 42     & 40    & 32    \\
permainan ke-5            & 24     & 16   \\ 
permainan ke-6            & \makebox[.15in][r]{8}          \\\hline
\end{tabular}
}
\caption{Penyelesaian Pascal untuk masalah pembagian taruhan.}
\label{table 3.11}
\end{table}               
\noindent Menurut ringkasan uraian Pascal, nilai permainan pertama sama dengan nilai
permainan kedua, sebagaimana dapat dibuktikan secara kombinatorial. Bilangan-bilangan pada
tiga baris pertama tabel meningkat, sedangkan bilangan pada baris keempat, kelima, dan
seterusnya menurun; Pascal mencatat bahwa pola ini tampak ganjil.\footnote{F.\ N. David,
op.\ cit., p.~235. Ringkasan editorial; ungkapan terjemahan modern tidak diterjemahkan
atau direproduksi.}
\par Mahasiswa dapat menyelidiki pertanyaan ini lebih lanjut menggunakan komputer dan
metode iterasi mundur Pascal untuk menghitung hasil harapan pada titik mana pun dalam
rangkaian.
\par Dalam risalahnya, Pascal memberikan bukti formal atas rumus kombinatorial Fermat,
serta bukti banyak sifat dasar lain dari bilangan binomial.  Banyak buktinya melibatkan
induksi dan termasuk di antara bukti-bukti pertama yang menggunakan metode ini.  Bukunya
menyatukan semua aspek bilangan dalam segitiga Pascal yang dikenal pada 1654; Edwards
menilai bahwa penamaan Segitiga Aritmetika menurut Pascal memiliki dasar sejarah yang
kuat.\index{PASCAL, B.|)}\index{FERMAT, P.|)}\footnote{A.~W.~F. Edwards, op.\ cit.,
p.~ix. Ringkasan editorial; ungkapan sumber tidak diterjemahkan atau direproduksi.}
\par Kajian serius pertama mengenai distribusi binomial dilakukan oleh James Bernoulli
dalam \emx{Ars Conjectandi} yang diterbitkan pada 1713.\index{BERNOULLI, J.}
\footnote{J. Bernoulli, \emx{Ars Conjectandi} (Basil: Thurnisiorum, 1713).}  Kita akan
kembali ke karya ini dalam catatan historis di Bab~\ref{chp 8}.

\exercises
\begin{LJSItem}

\i\label{exer 3.2.1} Hitunglah berikut ini:
\begin{enumerate}
\item ${6 \choose 3}$

\item $b(5,.2,4)$

\item ${7 \choose 2}$

\item ${{26} \choose {26}}$

\item $b(4,.2,3)$

\item ${6 \choose 2}$

\item ${{10} \choose 9}$

\item $b(8, .3, 5)$
\end{enumerate}

\i\label{exer 3.2.2} Ada berapa cara untuk memilih lima orang dari kelompok yang
terdiri atas sepuluh orang guna membentuk sebuah panitia?

\i\label{exer 3.2.3} Berapa banyak subhimpunan berunsur tujuh yang terdapat dalam
suatu himpunan berunsur sembilan?

\i\label{exer 3.2.4} Dengan menggunakan relasi dalam Persamaan~3.1, tulislah sebuah
program untuk menghitung segitiga Pascal dan menempatkan hasilnya dalam sebuah matriks.
Buatlah program Anda mencetak segitiga tersebut untuk $n = 10$.

\i\label{exer 3.2.5} Gunakan program {\bf BinomialProbabilities} untuk mencari peluang
bahwa, dalam 100 lemparan koin seimbang, banyaknya sisi kepala yang muncul terletak di
antara 35 dan 65, di antara 40 dan 60, serta di antara 45 dan 55.

\i\label{exer 3.2.6} Charles mengaku dapat membedakan bir dan ale dalam 75 persen
kesempatan.  Ruth bertaruh bahwa ia tidak mampu melakukannya dan sebenarnya hanya
menebak.  Untuk menyelesaikan persoalan ini, dibuatlah taruhan berikut: Charles diberi
sepuluh gelas kecil, yang masing-masing telah diisi dengan bir atau ale sebagaimana
ditentukan oleh lemparan koin seimbang.  Ia memenangkan taruhan jika menjawab benar
sebanyak tujuh kali atau lebih.  Carilah peluang Charles menang jika ia memang memiliki
kemampuan yang diakuinya.  Carilah peluang Ruth menang jika Charles hanya menebak.

\i\label{exer 3.2.7} Tunjukkan bahwa
$$ b(n,p,j) = \frac pq \left(\frac {n - j + 1}j \right) b(n,p,j - 1)\ ,
$$ untuk $j \ge 1$. Gunakan fakta ini untuk menentukan nilai atau nilai-nilai $j$ yang
membuat $b(n,p,j)$ mencapai nilai terbesarnya.  \emx {Petunjuk}: Perhatikan rasio-rasio
berturutan ketika $j$ bertambah.

\i\label{exer 3.2.8} Sebuah dadu dilempar 30 kali.  Berapakah peluang bahwa angka 6
muncul tepat 5 kali?  Berapa kali kemunculan angka 6 yang paling mungkin?

\i\label{exer 3.2.9} Carilah bilangan bulat $n$ dan $r$ sedemikian sehingga persamaan
berikut benar:
$$
 {13 \choose 5} + 2{13 \choose 6} + {13 \choose 7} = {n \choose r}\ .
$$

\i\label{exer 3.2.10} Dalam ujian benar-salah yang terdiri atas sepuluh soal, carilah
peluang seorang siswa memperoleh nilai 70 persen atau lebih hanya dengan menebak.
Jawablah pertanyaan yang sama jika ujian tersebut terdiri atas 30 soal dan jika ujian
tersebut terdiri atas 50 soal.

\i\label{exer 3.2.11} Sebuah restoran menawarkan pai apel dan pai bluberi serta
menyediakan kedua jenis pai dalam jumlah yang sama.  Setiap hari sepuluh pelanggan
memesan pai.  Dengan peluang yang sama, mereka memilih salah satu dari kedua jenis pai
itu.  Berapa potong dari setiap jenis pai yang harus disediakan pemilik agar peluang
setiap pelanggan memperoleh pai pilihannya kira-kira .95?

\i\label{exer 3.2.12} Satu tangan poker adalah himpunan 5 kartu yang dipilih secara
acak dari setumpuk 52 kartu.  Carilah peluang memperoleh

\begin{enumerate}
\item royal flush (sepuluh, jack, ratu, raja, dan as dari satu kembang).

\item straight flush (lima kartu berurutan dari satu kembang, tetapi bukan royal flush).

\item four of a kind (empat kartu dengan nilai muka yang sama).

\item full house (satu pasangan dan satu tripel, masing-masing bernilai muka sama).

\item flush (lima kartu dari satu kembang, tetapi bukan straight atau royal flush).

\item straight (lima kartu berurutan yang tidak semuanya dari kembang yang sama).
(Perhatikan bahwa dalam straight, as dapat dihitung sebagai kartu tinggi atau rendah.)
\end{enumerate}

\i\label{exer 3.2.13} Jika suatu himpunan memiliki $2n$ unsur, tunjukkan bahwa
himpunan itu memiliki lebih banyak subhimpunan berunsur $n$ daripada subhimpunan dengan
jumlah unsur lainnya.

\i\label{exer 3.2.14} Misalkan $b(2n,.5,n)$ adalah peluang bahwa dalam $2n$ lemparan
koin seimbang muncul tepat $n$ sisi kepala.  Dengan menggunakan rumus Stirling
(Teorema~\ref{thm 3.3}), tunjukkan bahwa $b(2n,.5,n)
\sim 1/\sqrt{\pi n}$.  Gunakan program {\bf BinomialProbabilities} untuk membandingkan
pendekatan ini dengan nilai eksaknya untuk $n = 10$ sampai~25.

\i\label{exer 3.2.15} Seorang pemain bisbol bernama Smith memiliki rata-rata pukulan
$.300$ dan biasanya mendapat giliran memukul tiga kali dalam satu pertandingan.
Anggaplah bahwa pukulan berhasil Smith dalam suatu pertandingan dapat dipandang sebagai
proses percobaan Bernoulli dengan peluang .3 untuk \emx{keberhasilan.} Carilah peluang
Smith memperoleh 0,~1,~2, dan~3 pukulan berhasil.

\i\label{exer 3.2.16} Tim sepak bola Universitas Siwash memainkan delapan pertandingan
dalam satu musim, dengan tiga kemenangan, tiga kekalahan, dan dua hasil seri.  Tunjukkan
bahwa banyaknya cara hal ini dapat terjadi adalah
$$ {8 \choose 3}{5 \choose 3} = \frac {8!}{3!\,3!\,2!}\ .
$$

\i\label{exer 3.2.17} Dengan menggunakan teknik pada Latihan~\ref{exer 3.2.16},
tunjukkan bahwa banyaknya cara menempatkan $n$ benda berbeda ke dalam tiga kotak,
dengan $a$ benda dalam kotak pertama, $b$ benda dalam kotak kedua, dan $c$ benda dalam
kotak ketiga, adalah $n!/(a!\,b!\,c!)$.

\i\label{exer 3.2.18} Baumgartner, Prosser, dan Crowell sedang memeriksa ujian
kalkulus.  Dalam ujian itu terdapat satu soal benar-salah yang terdiri atas sepuluh
bagian.  Baumgartner memperhatikan bahwa seorang siswa hanya menjawab benar dua dari
sepuluh bagian dan berkata, ``Siswa itu bahkan tidak cukup cerdas untuk melempar koin
dalam menentukan jawabannya."  ``Belum tentu," kata Prosser.  ``Dengan 340 siswa, saya
berani bertaruh bahwa jika mereka semua melempar koin untuk menentukan jawaban, akan ada
setidaknya satu lembar ujian dengan paling banyak dua jawaban benar." Crowell berkata,
``Saya sependapat dengan Prosser.  Bahkan, saya berani bertaruh bahwa semestinya ada
setidaknya satu lembar ujian tanpa satu pun jawaban benar jika semua orang hanya
menebak."  Siapakah yang benar dalam persoalan ini?

\i\label{exer 3.2.19} Satu tangan gin terdiri atas 10 kartu dari setumpuk 52 kartu.
Carilah peluang bahwa tangan gin tersebut memiliki

\begin{enumerate}
\item seluruh 10 kartu dari kembang yang sama.

\item tepat 4 kartu dari satu kembang dan masing-masing 3 kartu dari dua kembang lainnya.

\item sebaran kembang 4, 3, 2, 1.
\end{enumerate}

\i\label{exer 3.2.20} Satu tangan berisi enam kartu dibagikan dari setumpuk kartu
biasa.  Carilah peluang bahwa:
\begin{enumerate}
\item Keenam kartu tersebut adalah kartu hati.

\item Terdapat tiga as, dua raja, dan satu ratu.

\item Terdapat tiga kartu dari satu kembang dan tiga kartu dari kembang lainnya.
\end{enumerate}

\i\label{exer 3.2.21} Seorang perempuan ingin mewarnai kuku-kuku pada satu tangannya
dengan menggunakan paling banyak dua warna di antara merah, kuning, dan biru.  Ada
berapa cara ia dapat melakukannya?

\i\label{exer 3.2.22} Ada berapa cara untuk memasukkan enam surat yang tidak dapat
dibedakan ke dalam tiga kotak surat?  \emx {Petunjuk}: Salah satu cara merepresentasikan
hal ini diberikan oleh urutan $|$LL$|$L$|$LLL$|$, dengan tanda $|$ menyatakan sekat
antarkotak dan huruf L menyatakan surat.  Setiap cara yang mungkin dapat digambarkan
seperti itu.  Perhatikan bahwa kita memerlukan dua garis di kedua ujung, sedangkan dua
garis lainnya dan keenam huruf L dapat ditempatkan dalam urutan apa pun.

\i\label{exer 3.2.23} Dengan menggunakan metode dalam petunjuk
Latihan~\ref{exer 3.2.22}, tunjukkan bahwa $r$ benda yang tidak dapat dibedakan dapat
ditempatkan ke dalam $n$ kotak dengan
$$
 {{n + r - 1} \choose {n - 1}} = {{n + r - 1} \choose r}
$$ cara yang berbeda.

\i\label{exer 3.2.24} Sebuah biro perjalanan memperkirakan bahwa ketika 20 wisatawan
pergi ke sebuah tempat wisata yang memiliki sepuluh hotel, mereka menyebar seolah-olah
biro tersebut menempatkan 20 benda yang tidak dapat dibedakan ke dalam sepuluh kotak
yang dapat dibedakan.  Dengan menganggap model ini benar, carilah peluang tidak ada
hotel yang kosong ketika kelompok pertama yang terdiri atas 20 wisatawan tiba.

\i\label{exer 3.2.25} Sebuah lift\index{lift} membawa enam penumpang dan berhenti di
sepuluh lantai.  Kita dapat menetapkan dua ukuran peluang seragam yang berbeda untuk
cara para penumpang turun: (a)~kita menganggap para penumpang dapat dibedakan atau
(b)~kita menganggap mereka tidak dapat dibedakan (lihat Latihan~\ref{exer 3.2.23} untuk
kasus ini).  Untuk setiap kasus, hitunglah peluang bahwa semua penumpang turun di lantai
yang berbeda.

\i\label{exer 3.2.26} Anda bermain \emx{kepala atau ekor} dengan Prosser, tetapi Anda
mencurigai bahwa koinnya tidak seimbang.  Von~Neumann menyarankan prosedur berikut:
Lemparlah koin Prosser dua kali.  Jika hasilnya HT, sebut hasil tersebut \emx{menang;} jika
hasilnya TH, sebut hasil tersebut \emx{kalah.} Jika hasilnya TT atau HH, abaikan hasil itu
dan lempar kembali koin Prosser dua kali.  Lanjutkan sampai Anda memperoleh HT atau TH,
lalu sebut hasil tersebut menang atau kalah dalam satu permainan.  Ulangi prosedur ini
untuk setiap permainan.  Anggaplah koin Prosser menghasilkan sisi kepala dengan peluang
$p$.
\begin{enumerate}
\item Carilah peluang HT, TH, HH, dan TT dalam dua lemparan koin Prosser.

\item Dengan menggunakan bagian (a), tunjukkan bahwa peluang menang dalam setiap satu
permainan adalah 1/2, berapa pun nilai $p$.
\end{enumerate}

\i\label{exer 3.2.27} John mengaku memiliki kemampuan persepsi ekstrasensorik dan dapat
menentukan satu di antara dua simbol yang terdapat pada kartu tertelungkup (lihat
Contoh~\ref{exam 3.12}).  Untuk menguji kemampuannya, ia diminta melakukan hal tersebut
dalam serangkaian percobaan.  Misalkan hipotesis nolnya adalah bahwa ia sekadar menebak,
sehingga peluang ia menjawab benar setiap kali adalah 1/2, sedangkan hipotesis
alternatifnya adalah bahwa ia dapat menyebutkan simbol yang benar lebih dari separuh
waktu.  Rancanglah suatu uji yang peluang galat tipe 1-nya kurang dari .05 dan peluang
galat tipe 2-nya kurang dari .05 jika John dapat menyebutkan simbol yang benar dalam 75
persen kesempatan.

\i\label{exer 3.2.28} Dalam Contoh~\ref{exam 3.12}, anggaplah hipotesis alternatifnya
adalah $p = .8$ dan peluang setiap jenis galat diinginkan kurang dari .01.  Gunakan
program {\bf PowerCurve} untuk menentukan nilai $n$ dan $m$ yang memenuhi syarat ini.
Pilihlah $n$ sekecil mungkin.

\i\label{exer 3.2.29} Suatu obat dianggap efektif dengan peluang $p$ yang tidak
diketahui.  Untuk menaksir $p$, obat tersebut diberikan kepada $n$ pasien dan ternyata
efektif bagi $m$ pasien.  \emx{Metode kemungkinan maksimum}\index{Prinsip Kemungkinan\\
Maksimum} untuk menaksir $p$ menyatakan bahwa kita harus memilih nilai $p$ yang
memberikan peluang terbesar untuk memperoleh hasil yang kita peroleh dalam percobaan.
Dengan menganggap percobaan tersebut dapat dipandang sebagai proses percobaan Bernoulli
dengan peluang keberhasilan $p$, tunjukkan bahwa taksiran kemungkinan maksimum bagi $p$
adalah proporsi keberhasilan $m/n$.

\i\label{exer 3.2.30} Ingatlah bahwa dalam World Series, tim pertama yang memenangkan
empat pertandingan menjadi pemenang seri.  Seri tersebut berlangsung paling banyak
tujuh pertandingan.  Anggaplah Red Sox dan Mets sedang memainkan seri itu, dan Mets
memenangkan setiap pertandingan dengan peluang $p$.  Fermat mengamati bahwa meskipun
seri itu mungkin tidak berlangsung sampai tujuh pertandingan, peluang Mets memenangkan
seri sama dengan peluang mereka memenangkan empat pertandingan atau lebih dalam seri
yang dipaksakan berlangsung tujuh pertandingan, tanpa memandang siapa yang memenangkan
setiap pertandingan.

\begin{enumerate}
\item Dengan menggunakan program {\bf PowerCurve} pada Contoh~\ref{exam 3.12}, carilah
peluang Mets memenangkan seri untuk kasus $p = .5$, $p = .6$, dan $p =.7$.

\item Anggaplah Mets memiliki peluang .6 untuk memenangkan setiap pertandingan.  Gunakan
program {\bf PowerCurve} untuk mencari nilai $n$ sedemikian sehingga, jika seri
berlangsung sampai salah satu tim menjadi yang pertama memenangkan lebih dari separuh
pertandingan, Mets memiliki peluang 95 persen untuk memenangkan seri.  Pilihlah $n$
sekecil mungkin.
\end{enumerate}

\i\label{exer 3.2.31} Masing-masing dari empat mesin pada sebuah pesawat berfungsi
dengan baik dalam suatu penerbangan dengan peluang .99, dan mesin-mesin tersebut
berfungsi saling bebas.  Anggaplah pesawat dapat mendarat dengan selamat jika sedikitnya
dua mesinnya berfungsi dengan baik.  Berapakah peluang keadaan mesin-mesin tersebut
memungkinkan pesawat mendarat dengan selamat?

\i\label{exer 3.2.32} Seorang anak laki-laki tersesat ketika menuruni Gunung
Washington.  Pemimpin tim pencari memperkirakan bahwa anak itu menuruni sisi timur
dengan peluang $p$ dan sisi barat dengan peluang $1 - p$.  Tim pencarinya terdiri atas
$n$ orang yang akan mencari secara saling bebas; jika anak tersebut berada di sisi yang
sedang disusuri, setiap anggota akan menemukannya dengan peluang $u$.  Tentukan cara
pemimpin itu harus membagi $n$ orang tersebut menjadi dua kelompok untuk menyusuri kedua
sisi gunung agar peluang menemukan anak itu menjadi sebesar mungkin.  Bagaimana
jawabannya bergantung pada $u$?

\istar\label{exer 3.2.33} Sebanyak $2n$ bola dipilih secara acak dari kumpulan yang
terdiri atas $2n$ bola merah dan $2n$ bola biru.  Carilah ungkapan kombinatorial bagi
peluang bahwa bola-bola terpilih terbagi sama banyak menurut warna.  Gunakan rumus
Stirling untuk menaksir peluang ini.  Dengan menggunakan {\bf BinomialProbabilities},
bandingkan nilai eksaknya dengan pendekatan Stirling untuk $n = 20$.

\i\label{exer 3.2.34} Anggaplah bahwa setiap kali Anda membeli sekotak
Wheaties\index{Wheaties}, Anda memperoleh salah satu gambar dari $n$ pemain New York
Yankees\index{New York Yankees}.  Selama suatu periode, Anda membeli $m \geq n$ kotak
Wheaties.
\begin{enumerate}
\item Gunakan Teorema~\ref{thm 3.10} untuk menunjukkan bahwa peluang Anda memperoleh
seluruh $n$ gambar adalah
\begin{eqnarray*}
   1 &-& {n \choose 1} \left(\frac{n - 1}n\right)^m + {n \choose 2} \left(\frac{n - 2}n\right)^m -
           \cdots  \\
     &+& (-1)^{n - 1} {n \choose {n - 1}}\left(\frac 1n \right)^m.
\end{eqnarray*}
\emx {Petunjuk}: Misalkan $E_k$ adalah kejadian bahwa Anda tidak memperoleh gambar
pemain ke-$k$.

\item Tulislah sebuah program komputer untuk menghitung peluang ini.  Gunakan program
tersebut untuk mencari, bagi $n$ tertentu, nilai terkecil $m$ yang menghasilkan peluang
$\geq .5$ untuk memperoleh seluruh $n$ gambar.  Tinjau $n = 50$,~100, dan~150, lalu
tunjukkan bahwa $m = n\log n + n \log 2$ merupakan taksiran yang baik bagi banyaknya
kotak yang diperlukan.  (Untuk penurunan taksiran ini, lihat
Feller.\footnote{W. Feller, \emx{Introduction to Probability Theory and its Applications,}
vol.~I, 3rd~ed. (New York: John Wiley \& Sons, 1968), p.~106.})
\end{enumerate}

\istar\label{exer 3.2.35} Buktikan \emx{identitas binomial} berikut:
$$
 {2n \choose n} = \sum_{j = 0}^n { n \choose j}^2\ .
$$ \emx {Petunjuk}: Perhatikan sebuah guci yang berisi $n$ bola merah dan $n$ bola
biru.  Tunjukkan bahwa setiap ruas persamaan sama dengan banyaknya cara memilih $n$ bola
dari guci tersebut.

\i\label{exer 3.2.35.5} Misalkan $j$ dan $n$ bilangan bulat positif, dengan $j \le n$.
Suatu percobaan terdiri atas pemilihan secara acak sebuah tupel-$j$ bilangan bulat
\emx{positif} yang jumlahnya paling besar $n$.
\begin{enumerate}
\item Carilah ukuran ruang sampelnya.  \emx {Petunjuk}: Perhatikan $n$ bola yang tidak
dapat dibedakan dan disusun dalam satu baris.  Tempatkan $j$ penanda di antara
pasangan-pasangan bola yang bersebelahan, tanpa menempatkan dua penanda di antara
pasangan bola yang sama.  (Kita juga mengizinkan salah satu dari $n$ penanda ditempatkan
di ujung barisan bola.)  Tunjukkan bahwa terdapat korespondensi 1-1 antara
himpunan posisi penanda yang mungkin dan himpunan tupel-$j$ yang jumlahnya sedang kita
hitung.
\item Carilah peluang bahwa tupel-$j$ yang dipilih memuat sedikitnya satu angka 1.
\end{enumerate}

\i\label{exer 3.2.36} Misalkan $n\ (\mbox{mod}\ m)$ menyatakan sisa pembagian bilangan
bulat $n$ oleh bilangan bulat $m$.  Tulislah sebuah program komputer untuk menghitung
bilangan ${n \choose j}\ (\mbox{mod}\ m)$, dengan ${n \choose j}$ suatu koefisien
binomial dan $m$ suatu bilangan bulat.  Anda dapat melakukannya dengan menggunakan
relasi rekurensi untuk menghasilkan koefisien binomial dan mengerjakan seluruh
aritmetika memakai fungsi dasar mod($n,m$).  Usahakan agar program Anda membuat tabel
sebesar mungkin.  Jalankan program untuk kasus $m = 2$ sampai~7.  Apakah Anda melihat
suatu pola?  Khususnya, untuk kasus $m = 2$ dan $n$ merupakan pangkat~2, periksalah bahwa
semua entri pada baris ke-$(n - 1)$ adalah 1.  (Bilangan binomial yang bersesuaian
bersifat ganjil.)  Gunakan gambar Anda untuk menjelaskan mengapa hal ini benar.

\i\label{exer 3.2.37} Lucas\index{LUCAS, E.}\footnote{E. Lucas, ``Th\'eorie des Functions
Num\'eriques Simplement Periodiques," \emx{American J. Math.,} vol.~1 (1878), pp.~184-240,
289-321.} membuktikan hasil umum berikut yang berkaitan dengan Latihan~\ref{exer
3.2.36}.  Jika $p$ sebarang bilangan prima, maka ${n
\choose j}~ (\mbox{mod\ }p)$ dapat dicari sebagai berikut: Uraikan $n$ dan $j$ dalam
basis $p$, berturut-turut sebagai $n = s_0 + s_1p + s_2p^2 + \cdots + s_kp^k$ dan $j =
r_0 + r_1p + r_2p^2 + \cdots + r_kp^k$.  (Di sini $k$ dipilih cukup besar untuk
merepresentasikan semua bilangan dari 0 sampai $n$ dalam basis $p$ dengan menggunakan
$k$ digit.)  Misalkan $s = (s_0,s_1,s_2,\dots,s_k)$ dan $r =
(r_0,r_1,r_2,\dots,r_k)$.  Maka
$$
 {n \choose j}~(\mbox{mod\ }p) = \prod_{i = 0}^k   {{s_i} \choose {r_i}}~(\mbox{mod\
}p)\ .
$$ Sebagai contoh, jika $p = 7$, $n = 12$, dan $j = 9$, maka

\begin{eqnarray*} 12 & = & 5 \cdot 7^0 + 1 \cdot 7^1\ , \\
 9 & = & 2 \cdot 7^0 + 1 \cdot 7^1\ , 
\end{eqnarray*} sehingga
\begin{eqnarray*} s & = & (5, 1)\ , \\ r & = & (2, 1)\ , 
\end{eqnarray*} dan hasil ini menyatakan bahwa
$$ {12 \choose 9}~(\mbox{mod\ }p) = {5 \choose 2}  {1 \choose 1}~(\mbox{mod\ }7)\ .
$$ Karena ${12 \choose 9} = 220 = 3~(\mbox{mod\ }7)$ dan ${5 \choose 2} = 10 = 3~
(\mbox{mod\ }7)$, terlihat bahwa hasil tersebut benar untuk contoh ini.

Tunjukkan bahwa hasil ini menyiratkan bahwa, untuk $p = 2$, baris ke-$(p^k - 1)$ pada
segitiga Anda dalam Latihan~\ref{exer 3.2.36} tidak memiliki entri nol.

\i\label{exer 3.2.38} Buktikan bahwa peluang munculnya tepat $n$ sisi kepala dalam
$2n$ lemparan koin seimbang diberikan oleh hasil kali bilangan-bilangan ganjil sampai
$2n - 1$ dibagi dengan hasil kali bilangan-bilangan genap sampai $2n$.

\i\label{exer 3.2.39} Misalkan $n$ bilangan bulat positif dan anggaplah $j$ bilangan
bulat positif yang tidak melebihi $n/2$.  Tunjukkan bahwa dalam Teorema~\ref{thm 3.7},
jika perkalian dan pembagian dikerjakan secara berselang-seling, semua nilai antara
dalam perhitungan merupakan bilangan bulat.  Tunjukkan pula bahwa tidak satu pun nilai
antara tersebut melebihi nilai akhir.

\end{LJSItem}

\section{Pengocokan Kartu}\label{sec 3.3}\index{pengocokan kartu} 

Sebagian besar bagian ini didasarkan pada artikel karya Brad Mann,\footnote{B. Mann, ``How Many
Times Should You Shuffle a Deck of Cards?", \emx 
{UMAP Journal}, vol. 15, no. 4 (1994), pp. 303--331.}\index{MANN, B.}
yang merupakan uraian atas artikel karya David Bayer dan Persi Diaconis.\footnote{D. Bayer and
P. Diaconis, ``Trailing the Dovetail Shuffle to its Lair," \emx {Annals of Applied Probability},
vol. 2, no. 2 (1992), pp. 294--313.}\index{DIACONIS, P.}
\index{BAYER, D.}
\subsection*{Kocokan Riffle}
Diberikan setumpuk $n$ kartu, berapa kali kita harus mengocoknya agar menjadi ``acak"?
Tentu saja, jawabannya bergantung pada metode pengocokan yang digunakan dan pada makna
``acak" yang kita maksud.  Kita akan mulai mengkaji pertanyaan ini dengan meninjau model
standar bagi kocokan riffle.\index{kocokan riffle}
\par Kita mulai dengan setumpuk $n$ kartu, yang kita anggap diberi label bilangan bulat
dari 1 sampai $n$ dalam urutan menaik.  Kocokan riffle terdiri atas pemotongan\index{pemotongan}
tumpukan menjadi dua tumpukan dan penyisipan berselang-seling\index{penyisipan berselang-seling}
kedua tumpukan tersebut.  Sebagai contoh, jika $n = 6$, pengurutan awalnya adalah
$(1, 2, 3, 4, 5, 6)$, dan pemotongan dapat dilakukan di antara kartu 2 dan 3.  Hal ini
menghasilkan dua tumpukan, yaitu $(1, 2)$ dan $(3, 4, 5, 6)$.  Keduanya disisipkan secara
berselang-seling untuk membentuk pengurutan baru tumpukan kartu.  Sebagai contoh, kedua
tumpukan ini dapat membentuk pengurutan $(1, 3, 4, 2, 5, 6)$.  Untuk membahas pengocokan
semacam itu, kita perlu menetapkan suatu distribusi peluang pada himpunan semua pengocokan
yang mungkin.  Ada beberapa cara yang masuk akal untuk melakukannya.  Kita akan memberikan
beberapa strategi penetapan yang berbeda dan menunjukkan bahwa semuanya setara.  (Ini tidak
berarti bahwa penetapan tersebut merupakan satu-satunya penetapan yang masuk akal.)
Pertama, kita menetapkan peluang binomial $b(n, 1/2, k)$ pada kejadian bahwa pemotongan
dilakukan setelah kartu ke-$k$.  Selanjutnya, kita menganggap bahwa, untuk suatu
pemotongan, semua penyisipan berselang-seling yang mungkin memiliki peluang sama.  Jadi,
untuk melengkapi penetapan peluang, kita perlu menentukan banyaknya penyisipan
berselang-seling yang mungkin bagi dua tumpukan kartu yang masing-masing berisi $k$ dan
$n-k$ kartu.
\par Kita mulai dengan menuliskan tumpukan kedua dalam satu baris, dengan ruang di antara
setiap pasangan kartu yang berurutan serta ruang di awal dan akhir (jadi, terdapat
$n-k+1$ ruang).  Dengan pengembalian, kita memilih $k$ dari ruang-ruang ini dan
menempatkan kartu-kartu dari tumpukan pertama pada ruang yang dipilih.  Hal ini dapat
dilakukan dengan
$${{n}\choose{k}}$$ cara.  Jadi, peluang suatu penyisipan berselang-seling tertentu
seharusnya adalah
$${1\over{{n}\choose{k}}}\ .$$
\par Selanjutnya, perhatikan bahwa jika pengurutan baru bukan pengurutan identitas,
pengurutan itu merupakan hasil dari tepat satu pasangan pemotongan-penyisipan.  Jika
pengurutan baru adalah pengurutan identitas, pengurutan itu merupakan hasil dari salah
satu di antara $n+1$ pasangan pemotongan-penyisipan.
\par Kita mendefinisikan \emx{barisan naik}\index{barisan naik} dalam suatu pengurutan
sebagai subbarisan maksimal bilangan bulat berurutan dalam urutan menaik.  Sebagai contoh,
dalam pengurutan
$$(2, 3, 5, 1, 4, 7, 6)\ ,$$
terdapat 4 barisan naik, yaitu $(1)$, $(2, 3, 4)$, $(5, 6)$, dan $(7)$.  Mudah dilihat
bahwa suatu pengurutan merupakan hasil penerapan kocokan riffle pada pengurutan identitas
jika dan hanya jika pengurutan tersebut memiliki paling banyak dua barisan naik.  (Jika
pengurutan itu memiliki dua barisan naik, barisan-barisan tersebut bersesuaian dengan dua
tumpukan yang dihasilkan oleh pemotongan; jika pengurutan itu memiliki satu barisan naik,
pengurutan tersebut merupakan pengurutan identitas.)  Jadi, ruang sampel pengurutan yang
diperoleh dengan menerapkan kocokan riffle pada pengurutan identitas secara alami dapat
dideskripsikan sebagai himpunan semua pengurutan yang memiliki paling banyak dua barisan
naik.
\par Sekarang mudah untuk menetapkan distribusi peluang pada ruang sampel ini.  Setiap
pengurutan dengan dua barisan naik diberi nilai
$${{b(n, 1/2, k)}\over{{{n}\choose{k}}}} = {1\over{2^n}}\ ,$$ sedangkan pengurutan
identitas diberi nilai
$${{n+1}\over{2^n}}\ .$$
\par Ada cara lain untuk memandang kocokan riffle.  Kita dapat membayangkan memulai dengan
tumpukan kartu yang dipotong menjadi dua tumpukan seperti sebelumnya, dengan penetapan
peluang yang sama seperti sebelumnya, yaitu distribusi binomial.  Setelah memperoleh dua
tumpukan, kita mengambil kartu satu demi satu dari bagian bawah kedua tumpukan dan
menempatkannya pada satu tumpukan.  Jika pada suatu tahap proses ini kedua tumpukan
masing-masing berisi $k_1$ dan $k_2$ kartu, kita menganggap bahwa peluang kartu berikutnya
diambil dari suatu tumpukan sebanding dengan ukuran tumpukan tersebut saat itu.  Hal ini
menyiratkan bahwa peluang kartu berikutnya diambil dari tumpukan pertama sama dengan
$${{k_1}\over{k_1 + k_2}}\ ,$$
sedangkan peluang yang bersesuaian untuk tumpukan kedua adalah
$${{k_2}\over{k_1 + k_2}}\ .$$
Sekarang kita akan menunjukkan bahwa proses ini menetapkan peluang seragam pada setiap
penyisipan berselang-seling yang mungkin dari kedua tumpukan.
\par 
Sebagai contoh, misalkan suatu penyisipan berselang-seling terbentuk karena kartu-kartu
dari kedua tumpukan dipilih dalam suatu urutan.  Peluang hasil ini terjadi merupakan hasil
kali peluang pada setiap tahap proses, sebab pemilihan kartu pada setiap tahap dianggap
saling bebas dari pemilihan sebelumnya.  Setiap faktor dalam hasil kali ini berbentuk
$${{k_i}\over{k_1 + k_2}}\ ,$$ dengan $i = 1$ atau $2$, dan penyebut setiap faktor sama
dengan banyaknya kartu yang masih harus dipilih.  Jadi, penyebut peluang tersebut adalah
$n!$.  Pada saat sebuah kartu dipilih dari tumpukan yang berisi $i$ kartu, pembilang
faktor peluang yang bersesuaian adalah $i$, dan banyaknya kartu dalam tumpukan ini
berkurang sebanyak 1.  Jadi, pembilangnya adalah $k!(n-k)!$, sebab semua kartu dalam
kedua tumpukan pada akhirnya dipilih.  Oleh karena itu, proses ini menetapkan peluang
$${1\over{{n}\choose{k}}}$$ pada setiap penyisipan berselang-seling yang mungkin.
\par 
Sekarang kita beralih ke pertanyaan mengenai apa yang terjadi ketika kita melakukan
kocokan riffle sebanyak $s$ kali.  Jelas bahwa jika kita mulai dengan pengurutan
identitas, kita memperoleh pengurutan dengan paling banyak $2^s$ barisan naik, sebab satu
kocokan riffle menghasilkan paling banyak dua barisan naik dari setiap barisan naik dalam
pengurutan awal.  Bahkan, tidak sulit untuk melihat bahwa setiap pengurutan demikian
merupakan hasil dari $s$ kocokan riffle.  Dengan demikian, pertanyaannya menjadi: ada
berapa cara suatu pengurutan dengan $r$ barisan naik dapat terbentuk dengan menerapkan
$s$ kocokan riffle pada pengurutan identitas?  Untuk menjawab pertanyaan ini, kita beralih
ke gagasan kocokan-$a$.
\subsection*{Kocokan-$a$}
Ada beberapa cara untuk memvisualisasikan kocokan-$a$.  Salah satunya adalah membayangkan
makhluk bertangan $a$ yang diberi setumpuk kartu untuk dikocok dengan teknik riffle.
Makhluk itu secara alami memotong tumpukan menjadi $a$ tumpukan, lalu menyisipkan semuanya
secara berselang-seling.  (Bayangkan saja!)  Jadi, kocokan riffle biasa merupakan
kocokan-2.  Seperti dalam kocokan-2 biasa, kita mengizinkan beberapa tumpukan berisi 0
kartu.  Cara lain untuk memvisualisasikan kocokan-$a$ adalah dengan memikirkan inversnya,
yang disebut penguraian-$a$.  Gagasan ini dijelaskan dalam bukti teorema berikutnya.
\par Sekarang kita akan menunjukkan bahwa kocokan-$a$ yang diikuti kocokan-$b$ setara
dengan kocokan-$ab$.  Secara khusus, ini berarti bahwa $s$ kocokan riffle berturut-turut
setara dengan satu kocokan-$2^s$.  Kesetaraan ini dirumuskan secara tepat oleh teorema
berikut.
\noindent
\begin{theorem}\label{thm 3.3.1} Misalkan $a$ dan $b$ dua bilangan bulat positif.
Misalkan $S_{a,b}$ himpunan semua pasangan terurut dengan entri pertama berupa
kocokan-$a$ dan entri kedua berupa kocokan-$b$.  Misalkan $S_{ab}$ himpunan semua
kocokan-$ab$.  Maka terdapat korespondensi 1-1 antara $S_{a,b}$ dan $S_{ab}$ dengan sifat
berikut.  Misalkan $(T_1, T_2)$ bersesuaian dengan $T_3$.  Jika $T_1$ diterapkan pada
pengurutan identitas dan $T_2$ diterapkan pada pengurutan yang dihasilkan, pengurutan
akhirnya sama dengan pengurutan yang diperoleh dengan menerapkan $T_3$ pada pengurutan
identitas.
\proof Cara termudah untuk menjelaskan korespondensi yang diperlukan adalah melalui
gagasan penguraian\index{penguraian}.  Penguraian-$a$ dimulai dengan setumpuk $n$ kartu.
Satu demi satu, kartu diambil dari bagian atas tumpukan dan, dengan peluang yang sama,
ditempatkan di bagian bawah salah satu dari $a$ tumpukan, yang diberi label dari 0 sampai
$a-1$.  Setelah semua kartu dibagikan, kita menggabungkan tumpukan-tumpukan itu untuk
membentuk satu tumpukan dengan menempatkan tumpukan $i$ di atas tumpukan $i+1$, untuk
$0 \le i \le a-1$.  Mudah dilihat bahwa jika kita memulai dengan satu tumpukan kartu,
tepat ada satu cara untuk memotongnya sehingga diperoleh $a$ tumpukan yang dihasilkan
oleh penguraian-$a$; dengan $a$ tumpukan tersebut, tepat ada satu cara untuk menyisipkan
kartu-kartunya secara berselang-seling sehingga diperoleh tumpukan dalam urutan sebelum
penguraian dilakukan.  Jadi, penguraian-$a$ ini bersesuaian dengan tepat satu kocokan-$a$,
dan kocokan-$a$ tersebut merupakan invers dari penguraian-$a$ semula.
\par 
Jika kita menerapkan penguraian-$ab$ $U_3$ pada satu tumpukan kartu, kita memperoleh
himpunan yang terdiri atas $ab$ tumpukan, yang kemudian digabungkan menurut urutan untuk
membentuk satu tumpukan.  Kita melabeli tumpukan-tumpukan ini dengan pasangan terurut
bilangan bulat, dengan koordinat pertama berada di antara 0 dan $a-1$, sedangkan koordinat
kedua berada di antara 0 dan $b-1$.  Kemudian kita melabeli setiap kartu dengan label
tumpukannya.  Banyaknya label yang mungkin adalah $ab$, sebagaimana diperlukan.  Dengan
pelabelan ini, kita dapat menjelaskan cara mencari penguraian-$b$ dan penguraian-$a$
sedemikian sehingga, jika kedua penguraian tersebut diterapkan dalam urutan ini pada
tumpukan kartu, kita memperoleh himpunan $ab$ tumpukan yang sama dengan yang diperoleh
melalui penguraian-$ab$.
\par 
Untuk memperoleh penguraian-$b$ $U_2$, kita mengelompokkan tumpukan kartu menjadi $b$
tumpukan, dengan tumpukan ke-$i$ memuat semua kartu yang koordinat keduanya $i$, untuk
$0 \le i \le b-1$.  Kemudian tumpukan-tumpukan tersebut digabungkan untuk membentuk satu
tumpukan.  Penguraian-$a$ $U_1$ berlangsung dengan cara yang sama, kecuali bahwa koordinat
pertama pada label yang digunakan.  Selanjutnya, $a$ tumpukan yang dihasilkan digabungkan
untuk membentuk satu tumpukan.
\par Uraian di atas menunjukkan bahwa kartu-kartu yang akhirnya berada di atas adalah
semua kartu berlabel $(0, 0)$.  Kartu-kartu ini diikuti oleh kartu berlabel $(0, 1),\ (0, 2),$ $\  
\ldots,\ (0, b - 1),\ (1, 0),\ (1,1),\ldots,\ (a-1, b-1)$.  Selain itu,
urutan relatif setiap pasangan kartu yang memiliki label sama tidak pernah berubah.
Namun, hal ini tepat sama dengan penguraian-$ab$ jika, pada awal penguraian tersebut,
kita melabeli setiap kartu dengan salah satu label $(0, 0),\ (0, 1),\ \ldots,\ (0, b-1),\ (1, 0),\ (1,1),\ \ldots,\ (a-1, b-1)$.  Hal ini menyelesaikan bukti.
\end{theorem}
\par
Pada Gambar~\ref{fig 3.14}, kita menampilkan label untuk penguraian-2 atas setumpuk 10
kartu.  Terdapat 4 kartu berlabel 0 dan 6 kartu berlabel 1, sehingga jika penguraian-2
dilakukan, tumpukan pertama akan berisi 4 kartu dan tumpukan kedua akan berisi 6 kartu.
Ketika penguraian ini dilakukan, tumpukan kartu berakhir dalam pengurutan identitas.
\par
Pada Gambar~\ref{fig 3.15}, kita menampilkan label untuk penguraian-4 atas tumpukan kartu
yang sama (karena terdapat empat label yang digunakan).  Gambar ini juga dapat dipandang
sebagai contoh sepasang penguraian-2, sebagaimana dijelaskan dalam bukti di atas.
Penguraian-2 pertama akan menggunakan koordinat kedua pada label untuk menentukan
tumpukan.  Dalam kasus ini, kedua tumpukan memuat kartu-kartu yang nilainya adalah
$$
\{5, 1, 6, 2, 7\}\ {\rm dan}\ \{8, 9, 3, 4, 10\}\ .
$$ 
Setelah penguraian-2 ini dilakukan, tumpukan kartu berada dalam urutan yang ditunjukkan
pada Gambar~\ref{fig 3.14}, sebagaimana dapat diperiksa oleh pembaca.  Jika kita ingin
melakukan penguraian-4 pada tumpukan kartu dengan menggunakan label yang ditampilkan,
kita mengurutkan kartu secara leksikografis sehingga memperoleh empat tumpukan
$$
\{1, 2\},\ \{3, 4\},\ \{5, 6, 7\},\ {\rm dan}\ \{8, 9, 10\}\ .
$$
Ketika tumpukan-tumpukan ini digabungkan, kita sekali lagi memperoleh pengurutan identitas
tumpukan kartu.  Pokok teorema di atas adalah bahwa kedua prosedur pengurutan selalu
menghasilkan pengurutan awal yang sama.
\putfig{3truein}{PSfig3-14}{Sebelum penguraian-2.}{fig 3.14}  
\putfig{3truein}{PSfig3-15}{Sebelum penguraian-4.}{fig 3.15} 
\begin{theorem}\label{thm 3.3.2} Jika $D$ adalah sebarang pengurutan yang merupakan hasil
penerapan kocokan-$a$, kemudian kocokan-$b$, pada pengurutan identitas, maka peluang yang
ditetapkan pada $D$ oleh pasangan operasi ini sama dengan peluang yang ditetapkan pada
$D$ oleh proses penerapan kocokan-$ab$ pada pengurutan identitas.
\proof Sebut ruang sampel kocokan-$a$ sebagai $S_a$.  Jika kita melabeli tumpukan dengan
bilangan bulat dari $0$ sampai $a-1$, setiap pasangan pemotongan-penyisipan, yaitu
kocokan, bersesuaian dengan tepat satu bilangan bulat $n$ digit berbasis $a$, dengan
digit ke-$i$ pada bilangan bulat tersebut menyatakan tumpukan yang memuat
kartu ke-$i$.  Jadi, banyaknya pasangan pemotongan-penyisipan sama dengan banyaknya
bilangan bulat $n$ digit berbasis $a$, yaitu $a^n$.  Tentu saja, tidak
semua pasangan tersebut menghasilkan pengurutan yang berbeda.  Banyaknya pasangan yang
menghasilkan pengurutan tertentu akan dibahas kemudian.  Untuk keperluan kita, cukup
ditunjukkan bahwa pasangan pemotongan-penyisipanlah yang menentukan penetapan peluang.
\par Teorema sebelumnya menunjukkan bahwa terdapat korespondensi 1-1 antara $S_{a,b}$
dan $S_{ab}$.  Selain itu, unsur-unsur yang bersesuaian memberikan pengurutan yang sama
ketika diterapkan pada pengurutan identitas.  Diberikan sebarang pengurutan $D$, misalkan
$m_1$ adalah banyaknya unsur $S_{a,b}$ yang, ketika diterapkan pada pengurutan identitas,
menghasilkan $D$.  Misalkan $m_2$ adalah banyaknya unsur $S_{ab}$ yang, ketika diterapkan
pada pengurutan identitas, menghasilkan $D$.  Teorema sebelumnya menyiratkan bahwa $m_1
= m_2$.  Jadi, kedua himpunan menetapkan peluang
$${{m_1}\over{(ab)^n}}$$
pada $D$.  Hal ini menyelesaikan bukti.
\end{theorem}
\subsection*{Kaitan dengan Masalah Ulang Tahun} 
Ada hal lain yang dapat dikemukakan mengenai label yang diberikan kepada kartu oleh
penguraian-penguraian berturut-turut.  Misalkan kita melakukan penguraian-2 pada setumpuk $n$ kartu
sampai semua label pada kartu berbeda.  Mudah dilihat bahwa proses ini menghasilkan
setiap permutasi dengan peluang yang sama; artinya, proses ini merupakan proses acak.
Untuk melihatnya, perhatikan bahwa jika label-label menjadi berbeda pada penguraian-2
ke-$s$, urutan penguraian-2 ini dapat dipandang sebagai satu penguraian-$2^s$, dengan
setiap tumpukan yang ditentukan oleh penguraian tersebut memuat paling banyak satu kartu
(ingatlah bahwa tumpukan-tumpukan itu bersesuaian dengan label).  Jika setiap tumpukan
memuat paling banyak satu kartu, maka untuk sebarang dua kartu dalam tumpukan, peluang
kartu pertama memiliki label lebih rendah sama dengan peluangnya memiliki label lebih
tinggi daripada kartu kedua.  Jadi, setiap pengurutan yang mungkin memiliki peluang sama
untuk dihasilkan oleh penguraian-$2^s$ ini.
\par
Misalkan $T$ peubah acak yang menghitung banyaknya penguraian-2 sampai semua label berbeda.
Kita dapat memandang $T$ sebagai ukuran lamanya proses penguraian sampai keacakan tercapai.
Karena pengocokan dan penguraian merupakan proses yang saling invers, $T$ juga mengukur
banyaknya pengocokan yang diperlukan untuk mencapai keacakan.  Misalkan kita memiliki
setumpuk $n$ kartu dan ingin mencari $P(T \le s)$.  Nilai ini sama dengan $1 - P(T > s)$.
Namun, $T > s$ jika dan hanya jika setelah $s$ penguraian-2 tidak semua label berbeda.
Ini tidak lain dari masalah ulang tahun: kita mencari peluang sedikitnya dua orang
memiliki ulang tahun yang sama, jika terdapat $n$ orang dan $2^s$ kemungkinan tanggal
ulang tahun.  Dengan menggunakan rumus dari Contoh~\ref{exam 3.3}, kita memperoleh
\begin{equation}
P(T > s) = 1 - {{2^s} \choose n} \frac {n!}{2^{sn}}\ .
\label{eq 3.3.1}
\end{equation}
\par
Dalam Bab~\ref{chp 6}, kita akan mendefinisikan nilai rata-rata suatu peubah acak.  Dengan
menggunakan gagasan ini dan persamaan di atas, kita dapat menghitung nilai rata-rata
peubah acak $T$ (lihat Latihan~\ref{sec 6.1}.\ref{exer 6.1.42}).  Sebagai contoh, jika
$n = 52$, nilai rata-rata $T$ kira-kira 11.7.  Artinya, rata-rata diperlukan sekitar 12
kocokan riffle agar proses dapat dianggap acak.


\subsection*{Pasangan Pemotongan-Penyisipan dan Pengurutan}
Seperti dicatat dalam bukti Teorema~\ref{thm 3.3.2}, tidak semua pasangan
pemotongan-penyisipan menghasilkan pengurutan yang berbeda.  Namun, terdapat rumus
sederhana yang memberikan banyaknya pasangan semacam itu yang menghasilkan suatu
pengurutan tertentu.
\begin{theorem}\label{thm 3.3.3} Jika suatu pengurutan sepanjang $n$ memiliki $r$
barisan naik, banyaknya pasangan pemotongan-penyisipan dalam kocokan-$a$ atas pengurutan
identitas yang menghasilkan pengurutan tersebut adalah
$${{n + a - r}\choose{n}}\ .$$
\proof Untuk melihat mengapa hal ini benar, kita perlu menghitung banyaknya cara
pemotongan dalam kocokan-$a$ dapat dilakukan sehingga menghasilkan suatu pengurutan
dengan $r$ barisan naik.  Penyisipan berselang-seling dapat kita abaikan, sebab setelah
pemotongan dilakukan, paling banyak satu penyisipan akan menghasilkan suatu pengurutan
tertentu.  Karena pengurutan yang diberikan memiliki $r$ barisan naik, sebanyak $r-1$
titik pembagian dalam pemotongan telah ditentukan.  Sisanya, yaitu sebanyak $a - 1 -
(r - 1) = a - r$ titik pembagian, dapat ditempatkan di mana saja.  Banyaknya tempat untuk
menempatkan titik pembagian yang tersisa ini adalah $n+1$ (yaitu banyaknya ruang di antara
pasangan-pasangan kartu yang berurutan, termasuk posisi di awal dan akhir tumpukan).
Tempat-tempat ini dipilih dengan pengulangan yang diizinkan, sehingga banyaknya cara
melakukan pilihan tersebut adalah
$${{n + a - r}\choose{a-r}} = {{n + a - r}\choose{n}}\ .$$
\par Secara khusus, ini berarti bahwa jika $D$ merupakan pengurutan yang dihasilkan
dengan menerapkan kocokan-$a$ pada pengurutan identitas dan jika $D$ memiliki $r$ barisan
naik, peluang yang ditetapkan pada $D$ oleh proses ini adalah
$${{{{n + a - r}\choose{n}}}\over{a^n}}\ .$$ Hal ini menyelesaikan bukti.
\end{theorem}
\par
Teorema di atas menunjukkan bahwa informasi pokok mengenai peluang yang ditetapkan pada
suatu pengurutan dalam kocokan-$a$ hanyalah banyaknya barisan naik dalam pengurutan
tersebut.  Jadi, jika kita menentukan banyaknya pengurutan yang memuat tepat $r$ barisan
naik untuk setiap $r$ di antara 1 dan $n$, kita telah menentukan fungsi distribusi peubah
acak yang dihasilkan dengan menerapkan kocokan-$a$ acak pada pengurutan identitas.
\par Banyaknya pengurutan $\{1, 2, \ldots, n\}$ yang memiliki $r$ barisan naik
dinyatakan dengan $A(n, r)$ dan disebut bilangan Eulerian\index{bilangan Eulerian}.  Ada
banyak cara untuk menghitung nilai bilangan-bilangan ini; teorema berikut memberikan satu
metode rekursif yang langsung mengikuti apa yang telah kita ketahui mengenai kocokan-$a$.
\begin{theorem}\label{thm 3.3.4} Misalkan $a$ dan $n$ bilangan bulat positif.  Maka
\begin{equation}
a^n = \sum_{r = 1}^a {{n + a - r}\choose{n}}A(n, r)\ .
\label{eq 3.6}
\end{equation}
Jadi,
$$
A(n, a) = a^n - \sum_{r = 1}^{a-1}{{n + a - r}\choose{n}}A(n, r)\ .
$$
Selain itu,
$$A(n, 1) = 1\ .$$
\proof Persamaan kedua dapat digunakan untuk menghitung nilai bilangan Euler dan
langsung mengikuti Persamaan~\ref{eq 3.6}.  Persamaan terakhir merupakan akibat dari
fakta bahwa satu-satunya pengurutan $\{1, 2, \ldots, n\}$ yang memiliki satu barisan naik
adalah pengurutan identitas.  Jadi, tinggal membuktikan Persamaan~\ref{eq 3.6}.  Kita akan
menghitung himpunan kocokan-$a$ atas setumpuk $n$ kartu dengan dua cara.  Pertama, kita
mengetahui bahwa terdapat $a^n$ pengocokan demikian (hal ini dicatat dalam bukti
Teorema~\ref{thm 3.3.2}).  Namun, terdapat $A(n, r)$ pengurutan $\{1, 2, \ldots, n\}$
yang memiliki $r$ barisan naik, dan Teorema~\ref{thm 3.3.3} menyatakan bahwa untuk setiap
pengurutan demikian, tepat terdapat
$${{n+a-r}\choose n}$$ 
pasangan pemotongan-penyisipan yang menghasilkan pengurutan tersebut.  Oleh karena itu,
ruas kanan Persamaan~\ref{eq 3.6} menghitung himpunan kocokan-$a$ atas setumpuk $n$ kartu.
Hal ini menyelesaikan\linebreak[4] bukti.
\end{theorem}
\subsection*{Pengurutan Acak dan Proses Acak}
Sekarang kita beralih ke pertanyaan kedua yang diajukan pada awal bagian ini: Apa yang
kita maksud dengan pengurutan ``acak"\index{pengurutan, acak}\index{pengurutan acak}?
Agak menyesatkan jika suatu pengurutan tertentu dianggap acak atau tidak acak.  Jika kita
ingin memilih pengurutan acak dari himpunan semua pengurutan $\{1, 2, \ldots, n\}$, yang
kita maksud adalah bahwa setiap pengurutan harus dipilih dengan peluang yang sama; artinya,
setiap pengurutan sama ``acaknya" dengan pengurutan lainnya.
\par Kata ``acak" sebenarnya seharusnya digunakan untuk menjelaskan suatu proses.  Kita
akan mengatakan bahwa proses yang menghasilkan suatu objek dari himpunan objek (berhingga)
merupakan proses acak\index{proses, acak}\index{proses acak} jika setiap objek dalam
himpunan itu dihasilkan oleh proses tersebut dengan peluang yang sama.  Dalam keadaan
sekarang, objek-objeknya adalah pengurutan, dan proses yang menghasilkan objek-objek itu
adalah proses pengocokan.  Mudah dilihat bahwa tidak ada kocokan-$a$ yang benar-benar
merupakan proses acak, sebab jika $T_1$ dan $T_2$ adalah dua pengurutan dengan banyaknya
barisan naik yang berbeda, keduanya dihasilkan oleh kocokan-$a$ yang diterapkan pada
pengurutan identitas dengan peluang yang berbeda.
\subsection*{Jarak Variasi}
Alih-alih mensyaratkan bahwa suatu urutan pengocokan menghasilkan proses acak, kita akan
mendefinisikan ukuran yang menjelaskan seberapa jauh suatu proses tertentu dari proses
acak.  Misalkan $X$ sebarang proses yang menghasilkan pengurutan $\{1, 2, \ldots, n\}$.
Definisikan $f_X(\pi)$ sebagai peluang bahwa $X$ menghasilkan pengurutan $\pi$.
(Jadi, $X$ dapat dipandang sebagai peubah acak dengan fungsi distribusi $f$.)  Misalkan
$\Omega_n$ himpunan semua pengurutan $\{1, 2, \ldots, n\}$.  Terakhir, misalkan
$u(\pi) = 1/|\Omega_n|$ untuk semua $\pi \in \Omega_n$.  Fungsi $u$ merupakan fungsi
distribusi suatu proses acak yang menghasilkan pengurutan.  Untuk setiap pengurutan
$\pi \in \Omega_n$, kuantitas
$$|f_X(\pi) - u(\pi)|$$ adalah selisih antara peluang sebenarnya dan peluang yang
diinginkan bahwa $X$ menghasilkan $\pi$.  Jika kita menjumlahkannya untuk semua pengurutan
$\pi$ dan menyebut jumlah ini $S$, terlihat bahwa $S = 0$ jika dan hanya jika $X$ acak,
dan jika tidak, $S$ positif.  Mudah ditunjukkan bahwa nilai maksimum $S$ adalah 2, sehingga
kita mengalikan jumlah tersebut dengan $1/2$ agar nilainya berada dalam interval $[0, 1]$.
Jadi, kita memperoleh jumlah berikut sebagai rumus bagi \emx{jarak variasi}\index{jarak variasi}
antara kedua proses:
$$\parallel f_X - u \parallel = {1\over 2}\sum_{\pi \in \Omega_n} |f_X(\pi) - u(\pi)|\
.$$
\par Sekarang kita menerapkan gagasan ini pada kasus pengocokan.  Misalkan $X$ adalah
proses yang terdiri atas $s$ kocokan riffle berturut-turut yang diterapkan pada
pengurutan identitas.  Kita mengetahui bahwa $X$ juga dapat dipandang sebagai satu
kocokan-$2^s$.  Kita juga mengetahui bahwa $f_X$ konstan pada himpunan semua pengurutan
dengan $r$ barisan naik, dengan $r$ sebarang bilangan bulat positif.  Terakhir, kita
mengetahui nilai $f_X$ pada pengurutan dengan $r$ barisan naik serta banyaknya pengurutan
demikian.  Jadi, dalam kasus khusus ini, kita memiliki
$$\parallel f_X - u \parallel = {1\over 2}\sum_{r=1}^n A(n, r) 
\biggl|{{2^s + n - r}\choose{n}}/2^{ns} - {1\over{n!}}\biggr|\ .$$ Karena jumlah ini
hanya memiliki $n$ suku, mudah untuk menghitungnya bagi nilai $n$ berukuran sedang.
Untuk $n = 52$, kita memperoleh daftar nilai yang diberikan dalam
Tabel~\ref{table 3.12}.

\begin{table}
\centering
\begin{tabular}{ccrrrccllc}
\multicolumn{6}{c} {Banyaknya Kocokan Riffle} &
\multicolumn{4}{c} {Jarak Variasi}\\
\hline
&&&&1&&&&1\\
&&&&2&&&&1\\
&&&&3&&&&1\\
&&&&4&&&&0.9999995334\\
&&&&5&&&&0.9237329294\\
&&&&6&&&&0.6135495966\\
&&&&7&&&&0.3340609995\\
&&&&8&&&&0.1671586419\\
&&&&9&&&&0.0854201934\\
&&&&10&&&&0.0429455489\\
&&&&11&&&&0.0215023760\\
&&&&12&&&&0.0107548935\\
&&&&13&&&&0.0053779101\\
&&&&14&&&&0.0026890130
\end{tabular}
\caption{Jarak terhadap proses acak.}
\label{table 3.12}
\end{table}

\putfig{4truein}{PSfig3-13-5}{Jarak terhadap proses acak.}{fig 3.13}
 
\par Untuk membantu memahami data ini, data tersebut ditampilkan dalam bentuk grafik
pada Gambar~\ref{fig 3.13}.  Program {\bf VariationList}\index{VariationList (program)}
menghasilkan data yang ditampilkan dalam Tabel~\ref{table 3.12} dan
Gambar~\ref{fig 3.13}.  Terlihat bahwa sebelum terjadi 5 pengocokan, keluaran $X$ masih
sangat jauh dari acak.  Setelah 5 pengocokan, jarak dari proses acak pada dasarnya menjadi
setengah setiap kali pengocokan dilakukan.
\par Diberikan fungsi distribusi $f_X(\pi)$ dan $u(\pi)$ seperti di atas, ada cara lain
untuk memandang jarak variasi $\parallel f_X - u\parallel$.  Diberikan sebarang kejadian
$T$ (yang merupakan subhimpunan $S_n$), kita dapat menghitung peluangnya dalam proses $X$
dan dalam proses seragam.  Sebagai contoh, kita dapat membayangkan bahwa $T$ mewakili
himpunan semua permutasi ketika pemain pertama dalam permainan poker yang terdiri atas 7
pemain dibagikan straight flush (lima kartu berurutan dari kembang yang sama).  Menarik
untuk meninjau seberapa besar peluang kejadian ini setelah sejumlah pengocokan tertentu
berbeda dari peluang kejadian ini jika semua permutasi memiliki peluang sama.  Selisih
ini dapat dipandang sebagai penjelasan seberapa dekat proses $X$ dengan proses acak dalam
kaitannya dengan kejadian $T$.
\par Sekarang tinjau kejadian $T$ yang membuat nilai mutlak selisih kedua peluang tersebut
sebesar mungkin.  Dapat ditunjukkan bahwa nilai mutlak ini merupakan jarak variasi antara
proses $X$ dan proses seragam.  (Pembaca diminta membuktikan fakta ini dalam
Latihan~\ref{exer 3.3.4}.)
\par Kita baru saja melihat bahwa, untuk setumpuk 52 kartu, jarak variasi antara proses
7 kocokan riffle dan proses acak kira-kira $.334$.  Menarik untuk mencari suatu kejadian
$T$ sedemikian sehingga selisih antara peluang kedua proses menghasilkan $T$ mendekati
$.334$.  Kejadian dengan sifat ini dapat dijelaskan melalui permainan yang disebut
New-Age Solitaire.

\subsection*{New-Age Solitaire}\index{New-Age Solitaire}

Permainan ini diciptakan oleh Peter Doyle.  Permainan ini dimainkan dengan setumpuk
standar 52 kartu.  Kita membagikan kartu dalam keadaan terbuka, satu demi satu, ke
tumpukan buangan.  Jika menemukan kartu as, misalnya as Hati, kita menggunakannya untuk
memulai tumpukan Hati.  Setiap tumpukan kembang harus disusun secara berurutan dari as
sampai raja, hanya dengan menggunakan kartu yang dibagikan setelahnya.  Setelah semua
kartu dibagikan, kita mengambil tumpukan buangan dan melanjutkan permainan.  Kita
mendefinisikan kembang Yin\index{Yin} sebagai Hati dan Keriting, sedangkan kembang
Yang\index{Yang} sebagai Wajik dan Sekop.  Permainan berakhir ketika kedua tumpukan
kembang Yin telah selesai atau kedua tumpukan kembang Yang telah selesai.  Jelas bahwa
jika pengurutan tumpukan kartu dihasilkan oleh proses acak, peluang kedua tumpukan
kembang Yin selesai lebih dahulu tepat sama dengan 1/2.
\par Sekarang misalkan kita membeli setumpuk kartu baru, membuka segel kemasannya, dan
mengocok tumpukan itu dengan teknik riffle sebanyak 7 kali.  Jika percobaan ini dilakukan,
ternyata kembang Yin menang sekitar 75\% dari seluruh permainan.  Angka ini 25\% lebih
tinggi daripada yang kita peroleh jika tumpukan kartu benar-benar
diurutkan secara acak.  Penyimpangan ini cukup dekat dengan maksimum teoretis sebesar
$33.4$\% yang diperoleh di atas.
\par Mengapa kembang Yin begitu sering menang?  Dalam setumpuk kartu yang benar-benar
baru, kembang-kembangnya tersusun dalam urutan berikut dari atas ke bawah: as sampai raja
Hati, as sampai raja Keriting, raja sampai as Wajik, dan raja sampai as Sekop.  Perhatikan
bahwa jika kartu sama sekali tidak dikocok, tumpukan kembang Yin akan selesai dalam
putaran pertama, bahkan sebelum satu pun kartu kembang Yang terlihat.  Jika kita terus
memainkan permainan sampai tumpukan kembang Yang selesai, diperlukan 13 putaran melalui
tumpukan kartu.  Jadi, terlihat bahwa dalam tumpukan baru, kembang Yin berada dalam
urutan yang paling menguntungkan dan kembang Yang dalam urutan yang paling tidak
menguntungkan.  Setelah 7 kocokan riffle, keunggulan relatif kembang Yin atas kembang
Yang masih dipertahankan sampai batas tertentu.

\exercises
\begin{LJSItem}
\i\label{exer 3.3.1} Diberikan sebarang pengurutan $\sigma$ atas $\{1, 2, \ldots, n\}$,
kita dapat mendefinisikan $\sigma^{-1}$, yaitu pengurutan invers dari $\sigma$, sebagai
pengurutan yang unsur ke-$i$-nya merupakan posisi yang ditempati oleh $i$ dalam $\sigma$.
Sebagai contoh, jika $\sigma = (1, 3, 5, 2, 4, 7, 6)$, maka $\sigma^{-1} = (1, 4, 2, 5,
3, 7, 6)$.  (Jika pengurutan-pengurutan ini dipandang sebagai permutasi, maka
$\sigma^{-1}$ merupakan invers dari $\sigma$.)
\par Suatu \emx{penurunan}\index{penurunan} terjadi di antara dua posisi dalam suatu
pengurutan jika posisi kiri ditempati oleh bilangan yang lebih besar daripada bilangan
di posisi kanan.  Sebagai kesepakatan, kita mengatakan bahwa setiap pengurutan memiliki
penurunan setelah posisi terakhir.  Dalam contoh di atas, $\sigma^{-1}$ memiliki empat
penurunan.  Penurunan tersebut terjadi setelah posisi kedua, keempat, keenam, dan ketujuh.
Buktikan bahwa banyaknya barisan naik dalam suatu pengurutan $\sigma$ sama dengan banyaknya
penurunan dalam $\sigma^{-1}$.

\i\label{exer 3.3.2} Tunjukkan bahwa jika kita memulai dengan pengurutan identitas
$\{1, 2, \ldots, n\}$, peluang bahwa kocokan-$a$ menghasilkan pengurutan dengan tepat
$r$ barisan naik adalah
$${{{n + a - r}\choose{n}}\over{a^n}}A(n, r)\ ,$$ untuk $1 \le r \le a$.

\i\label{exer 3.3.3} Misalkan $D$ adalah setumpuk $n$ kartu.  Kita telah melihat bahwa terdapat
$a^n$ kocokan-$a$ atas $D$.  Pengodean himpunan penguraian-$a$ diberikan dalam bukti
Teorema~\ref{thm 3.3.1}.  Sekarang kita akan memberikan pengodean kocokan-$a$ yang
bersesuaian dengan pengodean penguraian-$a$.  Misalkan $S$ himpunan semua tupel-$n$
bilangan bulat, dengan setiap bilangan berada di antara 0 dan $a-1$.  Misalkan $M =
(m_1, m_2, \ldots, m_n)$ adalah sebarang unsur $S$.  Misalkan $n_i$ adalah banyaknya kemunculan
$i$ dalam $M$, untuk $0 \le i \le a-1$.  Anggaplah kita memulai dengan tumpukan kartu
dalam urutan menaik (yaitu, kartu-kartu diberi nomor dari 1 sampai $n$).  Kita memberi
$n_0$ kartu pertama label 0, $n_1$ kartu berikutnya label 1, dan seterusnya.
Kemudian, kocokan-$a$ yang bersesuaian dengan $M$ adalah pengocokan yang menghasilkan
pengurutan dengan kartu-kartu berlabel $i$ ditempatkan pada posisi dalam $M$ yang memuat
label $i$.  Kartu-kartu berlabel sama ditempatkan pada posisi-posisi tersebut dalam
urutan menaik berdasarkan nomornya.  Sebagai contoh, jika $n = 6$ dan $a = 3$, misalkan
$M = (1, 0, 2, 2, 0, 2)$.  Maka $n_0 = 2,\ n_1 = 1,$ dan $n_2 = 3$.  Jadi, kita memberi
kartu 1 dan 2 label 0, kartu 3 label 1, serta kartu 4, 5, dan 6 label 2.
Kemudian kartu 1 dan 2 ditempatkan pada posisi 2 dan 5, kartu 3 ditempatkan pada posisi 1,
serta kartu 4, 5, dan 6 ditempatkan pada posisi 3, 4, dan 6, sehingga menghasilkan
pengurutan $(3, 1, 4, 5, 2, 6)$.
\begin{enumerate}
\item Dengan menggunakan pengodean ini, tunjukkan bahwa dalam kocokan-$a$, peluang kartu
pertama (yaitu, kartu nomor 1) berpindah ke posisi ke-$i$ diberikan oleh ungkapan berikut:
$$ {{(a-1)^{i-1}a^{n-i} + (a-2)^{i-1}(a-1)^{n-i} + \cdots + 1^{i-1}2^{n-i}}\over{a^n}}\ .
$$
\item Berikan taksiran akurat bagi peluang bahwa, dalam tiga kocokan riffle atas setumpuk
52 kartu, kartu pertama berakhir di salah satu dari 26 posisi pertama.  Dengan menggunakan
komputer, taksirlah secara akurat peluang kejadian yang sama setelah tujuh kocokan riffle.
\end{enumerate}

\i\label{exer 3.3.4} Misalkan $X$ menyatakan suatu proses tertentu yang menghasilkan
unsur-unsur $S_n$, dan misalkan $U$ menyatakan proses seragam.  Misalkan fungsi distribusi
kedua proses ini berturut-turut dinyatakan dengan $f_X$ dan $u$.  Tunjukkan bahwa jarak
variasi
\newline	
$\parallel f_X - u\parallel$ sama dengan
$$\max_{T \subset S_n} \sum_{\pi \in T} \Bigl(f_X(\pi) - u(\pi)\Bigr)\ .$$ 
\emx {Petunjuk}:
Tuliskan permutasi-permutasi dalam $S_n$ menurut urutan menurun dari selisih $f_X(\pi) -
u(\pi)$.

\i\label{exer 3.3.5} Tinjaulah proses yang dijelaskan dalam teks, yaitu setumpuk $n$
kartu berulang kali diberi label dan dikenai penguraian-2 dengan cara yang dijelaskan
dalam bukti Teorema~\ref{thm 3.3.1}.  (Lihat Gambar~\ref{fig 3.12} dan
\ref{fig 3.13}.)  Proses berlanjut sampai semua label berbeda.  Tunjukkan bahwa proses
tidak pernah berakhir sebelum sedikitnya $\lceil \log_2(n) \rceil$ penguraian dilakukan.
\end{LJSItem}
